\exercisesection{}

In the following exercises, when E is a Banach space, we denote by \(\mathcal{C}alk(E)\) the Calkin algebra \(\mathcal{L}(E)/\mathcal{L}^{c}(E)\) . We denote by \(\pi\) the canonical projection from \(\mathcal{L}(E)\) onto \(\mathcal{C}alk(E)\) .

If E is a complex Hilbert space, we equip \(\mathcal{C}alk(\mathrm{E})\) with the involution induced from the involution of \(\mathcal{L}(\mathrm{E})\) by passage to the quotient.

\begin{enumerate}
\def\labelenumi{\arabic{enumi})}
\tightlist
\item
  Let E be a locally convex space.
\end{enumerate}

\begin{enumerate}
\def\labelenumi{\alph{enumi})}
\item
  Suppose that every neighborhood of 0 in E contains a vector subspace of E of finite codimension. Show that the space \(\mathcal{F}(\mathrm{E})\) of Fredholm endomorphisms of E is not closed in \(\mathcal{L}(\mathrm{E})\) endowed with the topology of bounded convergence.
\item
  Give an example of an infinite-dimensional Fréchet space such that every neighborhood of 0 contains a vector subspace of E of finite codimension. c) Let \(E_{\sigma}\) be an infinite-dimensional Banach space equipped with the weak topology. The set \(\mathcal{F}(E_{\sigma})\) is not open in \(\mathcal{L}(E_{\sigma})\) .
\end{enumerate}

\begin{enumerate}
\def\labelenumi{\arabic{enumi})}
\setcounter{enumi}{1}
\tightlist
\item
  Let E be an infinite-dimensional Banach space and \(u \in \mathcal{L}(\mathrm{E})\) . We denote by \(\mathrm{S}(u)\) the set of complex numbers \(\lambda\) such that there exist a real number \(\delta > 0\) and a finite-codimensional subspace F of E such that \(\| u(x) - \lambda x \| \geqslant \delta \| x \|\) for all \(x \in \mathrm{F}\) .
\end{enumerate}

\begin{enumerate}
\def\labelenumi{\alph{enumi})}
\item
  The set \(S(u)\) is open in C and contains the complement in C of the spectrum of u.
\item
  If \(\lambda \in \mathrm{S}(u)\) , the endomorphism \(u - \lambda 1_{\mathrm{E}}\) has a closed image and a finite-dimensional kernel.
\item
  The complement \(\mathrm{F}(u) = \mathbf{C} - \mathrm{S}(u)\) is not empty. (Show that an element of \(\mathrm{Sp}_{\mathrm{e}}(u)\) of maximal modulus belongs to \(\mathrm{F}(u)\) .)
\end{enumerate}

\begin{enumerate}
\def\labelenumi{\arabic{enumi})}
\setcounter{enumi}{2}
\tightlist
\item
  Let E be an infinite-dimensional complex Hilbert space.
\end{enumerate}

\begin{enumerate}
\def\labelenumi{\alph{enumi})}
\item
  The involutive algebra \(\mathcal{C}alk(E)\) is a star algebra.
\item
  Let u be an element of \(\mathcal{L}(\mathrm{E})\) such that \(u^{*}u\) is compact. Show that u is compact.
\item
  Let G be the group of invertible elements in \(\mathcal{C}alk(\mathrm{E})\) and \(G_{0}\) its neutral component. Let \(p: G \to G/G_{0}\) be the canonical projection. There exists a unique group isomorphism \(\alpha: Z \to G/G_{0}\) such that \(\alpha(\operatorname{ind}(u)) = p(\pi(u))\) for every Fredholm endomorphism u of E.
\end{enumerate}

\begin{enumerate}
\def\labelenumi{\arabic{enumi})}
\setcounter{enumi}{3}
\tightlist
\item
  Let E be a complex Hilbert space. For every invertible element \(\xi\) of \(\mathcal{C}alk(E)\) , we call the index of \(\xi\) the index of any Fredholm endomorphism \(u\) of E such that \(\pi(u) = \xi\) .
\end{enumerate}

\begin{enumerate}
\def\labelenumi{\alph{enumi})}
\item
  Show that for every \(\xi \in \mathcal{C}alk(\mathrm{E})\) that is Hermitian, there exists \(u \in \mathcal{L}(\mathrm{E})\) that is Hermitian such that \(\pi(u) = \xi\) .
\item
  Show that for every \(\xi \in \mathcal{C}alk(\mathrm{E})\) that is unitary and has index zero, there exists \(u \in \mathcal{L}(\mathrm{E})\) that is unitary such that \(\pi(u) = \xi\) . (Let \(v\) be invertible such that \(\pi(v) = \xi\) ; consider the polar decomposition \(v = j|v|\) of \(v\) and verify that \(|v| - 1_{\mathrm{E}}\) is compact, then deduce that \(v - j\) is compact.)
\item
  Show that for every \(\xi \in \mathcal{C}alk(\mathrm{E})\) that is unitary such that \(\xi^2 = 1\) , there exists \(u \in \mathcal{L}(\mathrm{E})\) that is unitary such that \(u^2 = 1\) and \(\pi(u) = \xi\) .
\end{enumerate}

\(d)\) For every element \(p \in \mathcal{C}alk(\mathrm{E})\) such that \(p^2 = p\) , there exists a projector \(\mathrm{P} \in \mathcal{L}(\mathrm{E})\) such that \(\pi(\mathrm{P}) = p\) .

\begin{enumerate}
\def\labelenumi{\arabic{enumi})}
\setcounter{enumi}{4}
\tightlist
\item
  Let E be a complex Hilbert space.
\end{enumerate}

\begin{enumerate}
\def\labelenumi{\alph{enumi})}
\tightlist
\item
  Let \(n\) be an integer with \(n \geqslant 2\) . We set \(\mathrm{I} = \mathbf{N} \times \{0, \ldots, n\}\) . Let \((\sigma_i)_{1 \leqslant i \leqslant n}\) be the family of permutations of I defined by
\end{enumerate}

\[
\begin{array}{c} \sigma_ {i} (m, i) = (m - 1, i) \text {for} m \geqslant 1 \\ \sigma_ {i} (0, i) = (0, i - 1) \\ \sigma_ {i} (m, i - 1) = (m + 1, i - 1) \text {for} m \geqslant 0 \\ \sigma_ {i} (m, j) = (m, j) \text {for} m \geqslant 0 \text {and} j \notin \{i, i - 1 \}. \end{array}
\]

\begin{enumerate}
\def\labelenumi{\alph{enumi})}
\setcounter{enumi}{1}
\item
  We denote by \((e_{m,i})\) the canonical orthonormal basis of \(\ell^{2}(\mathrm{I})\) . Let \(u_{i}\) be the endomorphism of \(\ell^{2}(\mathrm{I})\) such that \(u_{i}(e_{m,j}) = e_{\sigma_{i}(m,j)}\) for all \((m,j) \in \mathrm{I}\) . The endomorphism \(u_{i}\) is unitary.
\item
  The elements \(\pi(u_i)\) are pairwise commuting.
\end{enumerate}

\(d)\) There does not exist a family \((v_{1},\ldots,v_{n})\) of unitary endomorphisms of E, pairwise commuting, such that \(\pi(v_{i})=\pi(u_{i})\) for all \(i\) .

\begin{enumerate}
\def\labelenumi{\arabic{enumi})}
\setcounter{enumi}{5}
\tightlist
\item
  Let E be a complex Hilbert space. Recall that \(\mathcal{C}alk(E)\) is a star algebra (exercise 3).
\end{enumerate}

  Let \(\mathfrak{F}\) be a nonprincipal ultrafilter on \(\mathbf{N}\) . Let F be the vector space of bounded sequences in E that converge weakly to 0 along \(\mathfrak{F}\) and let \(F_0\) be the subspace of F consisting of sequences that converge in norm to 0 along \(\mathfrak{F}\) . We set \(\widetilde{\mathrm{H}} = \mathrm{F} / \mathrm{F}_0\) . a) For all elements \(x = (x_n)\) and \(y = (y_n)\) of F, show that the sequence \((\langle x_n | y_n \rangle)_{n \in \mathbf{N}}\) converges along \(\mathfrak{F}\) . We denote by \(\langle x | y \rangle_{\mathrm{F}}\) its limit.

\begin{enumerate}
\def\labelenumi{\alph{enumi})}
\setcounter{enumi}{1}
\item
  Show that \((x,y)\mapsto\langle x\mid y\rangle_{\mathrm{F}}\) is a positive Hermitian form on F, and that \(\langle x\mid x\rangle_{F}=0\) if and only if \(x\in F_{0}\) ; this Hermitian form therefore defines a separated pre-Hilbert space structure on \(\widetilde{H}\).
\item
  Let H be the Hilbert space obtained by completing the separated pre-Hilbert space \(\widetilde{H}\) . Show that for every u in \(\mathcal{L}(E)\) , one obtains an element \(\varrho(u)\) of \(\mathcal{L}(H)\) by passing to the quotient of the linear map \((x_{n})\mapsto(u(x_{n}))\) from F to F. d) The map \(u\mapsto\varrho(u)\) induces an isometric morphism of C*-algebras \(\mathcal{C}alk(E)\to\mathcal{L}(H)\) .
\end{enumerate}

(We will see later that every C*-algebra A has an isometric representation \(A \to \mathcal{L}(E)\) for a certain Hilbert space E, cf. th. 2 of V, p. 442.)

\begin{enumerate}
\def\labelenumi{\arabic{enumi})}
\setcounter{enumi}{6}
\tightlist
\item
  We equip the unit circle \(S_{1} \subset \mathbf{C}\) with the normalized Haar measure. Let E be the complex Hilbert space \(\mathrm{L}^{2}(\mathbf{S}_{1})\) . For \(n \in \mathbf{Z}\) , let \(e_{n}\) be the element \(z \mapsto z^{n}\) of E. Let F ⊂ E be the closed subspace generated by the \(e_{n}\) for \(n \geqslant 0\) and let \(p: E \to F\) be the orthogonal projector. For \(f \in L^{\infty}(\mathbf{S}_{1})\) we define linear maps \(M_{f}\) (resp. \(T_{f}\) ) on E (resp. F) by \(\mathrm{M}_{f}(g) = fg\) for \(g \in E\) and \(T_{f} = p \circ M_{f}\) . One says that \(T_{f}\) is the Toeplitz operator with symbol f. (a) The linear maps \(M_{f}\) and \(T_{f}\) are continuous. The spectrum of \(M_{f}\) is the closure of the set of Fourier coefficients \(c_{n} = \langle f | e_{n} \rangle\) of f. The endomorphism \(M_{f}\) is a Fredholm endomorphism if and only if it is a Riesz endomorphism.
\end{enumerate}

\begin{enumerate}
\def\labelenumi{\alph{enumi})}
\setcounter{enumi}{1}
\item
  The endomorphisms \(T_{e_{n}}\) for \(n \in Z\) are Fredholm endomorphisms and \(\text{ind}(\mathrm{T}_{e_{n}}) = -n\) for all \(n \in Z\) .
\item
  One has \(\mathrm{T}_{\overline{f}} = \mathrm{T}_f^*\) for every function \(f\in \mathrm{L}^{\infty}(\mathbf{S}_1)\) .
\item
  If \(T_{f}\) is an isomorphism, then \(M_{f}\) is an isomorphism and f is invertible in \(\mathrm{L}^{\infty}(\mathbf{S}_{1})\) . Deduce from this that the spectrum of \(M_{f}\) is contained in the spectrum of \(T_{f}\) and that \(\|T_{f}\|=\|f\|\) .
\item
  Let \(f \in \mathcal{C}(\mathbf{S}_1)\) be a continuous function and \(g \in \mathrm{L}^{\infty}(\mathbf{S}_1)\) . The endomorphisms \(\mathrm{T}_f\mathrm{T}_g - \mathrm{T}_{fg}\) and \(\mathrm{T}_g\mathrm{T}_f - \mathrm{T}_{gf}\) are compact. (First consider the case where \(f = e_n\) .)
\item
  Let \(\mathcal{C}alk(\mathrm{F})\) be the Calkin algebra of F and \(\pi\colon\mathcal{L}(\mathrm{F})\to\mathcal{C}alk(\mathrm{F})\) be the canonical projection. Show that \(f\mapsto\pi(\mathrm{T}_{f})\) is an injective morphism of involutive algebras from \(\mathcal{C}(\mathbf{S}_{1})\) into \(\mathcal{C}alk(\mathrm{F})\) (To prove injectivity, one may use prop. 1 of I, p. 31.)
\end{enumerate}

\(g)\) Let \(f \in \mathcal{C}(\mathbf{S}_1)\) . Then \(\mathrm{T}_f\) is a Fredholm endomorphism if and only if \(f\) does not vanish on \(\mathbf{S}_1\) . (Use prop. 5 of I, p. 106.)

\(h)\) Let \(f \in \mathcal{C}(\mathbf{S}_1)\) such that \(f\) does not vanish on \(\mathbf{S}_1\) . The index of \(\mathrm{T}_f\) is equal to \(n \in \mathbf{Z}\) if and only if the path \(t \mapsto f(e^{2i\pi t})\) in \(\mathbf{C} - \{0\}\) is homotopic to the path \(t \mapsto e_n(t)\) (cf. TA, III, p. 229, def. 1).

\begin{enumerate}
\def\labelenumi{\arabic{enumi})}
\setcounter{enumi}{7}
\tightlist
\item
  We resume the notation of the preceding exercise.
\end{enumerate}

\begin{enumerate}
\def\labelenumi{\alph{enumi})}
\item
  We denote by \(\mathrm{T}(\mathbf{S}_{1})\) the closed subalgebra of \(\mathcal{L}(\mathrm{F})\) generated by the Toeplitz operators \(T_{e_{1}}\) and \(T_{e_{-1}}\) . Show that \(\mathrm{T}(\mathbf{S}_{1})\) is a star subalgebra of \(\mathcal{L}(\mathrm{F})\) .
\item
  Let \(f \in \mathcal{C}(\mathbf{S}_1)\) . Show that \(T_f\) belongs to \(T(\mathbf{S}_1)\) .
\item
  The representation of \(\mathrm{T}(\mathbf{S}_{1})\) on F is irreducible; the star algebra \(\mathrm{T}(\mathbf{S}_{1})\) contains \(\mathcal{L}^{\mathrm{c}}(\mathrm{F})\) .
\end{enumerate}

\(d)\) One has \(u \in \mathrm{T}(\mathbf{S}_1)\) if and only if there exists \(f \in \mathcal{C}(\mathbf{S}_1)\) and \(h \in \mathcal{L}^c(\mathrm{F})\) such that \(u = \mathrm{T}_f + h\) . (Show that the set of endomorphisms of the form \(\mathrm{T}_f + h\) forms a closed star algebra in \(\mathcal{L}(\mathrm{F})\) .)

\begin{enumerate}
\def\labelenumi{\alph{enumi})}
\setcounter{enumi}{4}
\tightlist
\item
  There exists an exact sequence of morphisms of involutive algebras
\end{enumerate}

\[
0 \longrightarrow \mathcal {L} ^ {\mathrm{c}} (\mathrm{F}) \longrightarrow \mathrm{T} (\mathbf {S} _ {1}) \stackrel {\varrho} {\longrightarrow} \mathcal {C} (\mathbf {S} _ {1}) \longrightarrow 0
\]

such that \(\varrho(\mathrm{T}_f) = f\) for every function \(f \in \mathcal{C}(\mathbf{S}_1)\) .

\begin{enumerate}
\def\labelenumi{\arabic{enumi})}
\setcounter{enumi}{8}
\item
  Let B be an infinite-dimensional Banach space. We denote by E the strong dual of B and by F the dual of B equipped with the topology \(\sigma(\mathrm{B}^{\prime},\mathrm{B})\) . Show that there exists a continuous bijective linear map u from E into F and a compact linear map h from E into F such that \(u + h\) has a cokernel of infinite dimension.
\item
  Let E be a separated locally convex space, and let u and h be commuting elements of \(\mathcal{L}(\mathrm{E})\) . Suppose there exists an integer n such that \(h^{n}\) is a compact endomorphism.
\end{enumerate}

\begin{enumerate}
\def\labelenumi{\alph{enumi})}
\item
  If u is a Fredholm endomorphism, then u-h is a Fredholm endomorphism.
\item
  If \(u\) is a Riesz endomorphism, then \(u - h\) is a Riesz endomorphism.
\end{enumerate}

\begin{enumerate}
\def\labelenumi{\arabic{enumi})}
\setcounter{enumi}{10}
\tightlist
\item
  Let \(\mathrm{E} = \ell^2 (\mathbf{N})\) and let \((e_n)\) be the canonical basis of E. We define \(u\in \mathcal{L}(\mathrm{E})\) by \(u(e_i) = e_{i + 1}\) for \(i\in \mathbf{N}\) . Let \(\mathrm{F} = \mathrm{E}\oplus \mathrm{E}\) and \(v = u\oplus (u^{*} + 2\cdot 1_{\mathrm{E}})\in \mathcal{L}(\mathrm{F})\) . Let \(\overline{v} = \pi (v)\) be the image of \(v\) in the Calkin algebra \(\mathcal{Calk}(\mathrm{F})\) .
\end{enumerate}

\begin{enumerate}
\def\labelenumi{\alph{enumi})}
\item
  The endomorphism \(v\) is a Fredholm endomorphism of index -1, and \(v - 2 \cdot 1_{\mathrm{E}}\) is a Fredholm endomorphism of index 1.
\item
  Let \(p = \mathrm{X}(\mathrm{X} - 2) \in \mathbf{C}[\mathrm{X}]\) and let S be the exponential spectrum of \(\overline{v}\) (exercise 5 of I, p. 173). We have \(0 \in p(\mathrm{S})\) but 0 does not belong to the exponential spectrum of \(p(\overline{v})\) (Use Exercise 3.)
\end{enumerate}

\exercisesection{}

\begin{enumerate}
\def\labelenumi{\arabic{enumi})}
\tightlist
\item
  Let A be a compact subset of \(\mathbf{C}\) which is equal to the closure of its interior V. We set \(\mathcal{C}(\mathrm{A}) = \mathcal{C}(\mathrm{A};\mathbf{C})\) . Let \(B(A)\) be the vector subspace of \(\mathcal{C}(\mathrm{A})\) consisting of functions \(f\in \mathcal{C}(\mathrm{A})\) whose restriction to V is holomorphic;
\end{enumerate}

it is a Banach subalgebra of \(\mathcal{C}(\mathrm{A})\) . Let \(u \in \mathcal{L}(\mathrm{B}(\mathrm{A}))\) be the linear map such that \(u(f)(z) = zf(z)\) for all \(z \in \mathrm{A}\) and every \(f \in \mathrm{B}(\mathrm{A})\) .

\begin{enumerate}
\def\labelenumi{\alph{enumi})}
\item
  The spectrum of \(u\) is equal to A.
\item
  For every \(z \in V\) the endomorphism u - z of B(A) is a Fredholm endomorphism of index 1.
\item
  For every \(z \in A - V\) the endomorphism u - z of B(A) is not a Fredholm endomorphism.
\item
  The stable spectrum of u is the boundary of A, the sensitive spectrum of u is empty, and the essential spectrum of u is A.
\item
  We assume that the complement of A in \(\mathbf{C}\) is connected. For every \(h \in \mathcal{L}^{c}(\mathrm{B}(\mathrm{A}))\) , the essential spectrum of \(u + h\) is equal to A.
\end{enumerate}

\begin{enumerate}
\def\labelenumi{\arabic{enumi})}
\setcounter{enumi}{1}
\tightlist
\item
  We keep the notation from the previous exercise. We assume that A contains the unit circle \(U\) in \(\mathbf{C}\) and that \(0 \notin A\) . We denote by \(\varphi_{A}\) the characteristic function of A.
\end{enumerate}

\begin{enumerate}
\def\labelenumi{\alph{enumi})}
\tightlist
\item
  The map
\end{enumerate}

\[
\ell \colon f \mapsto \frac {1}{2 i \pi} \int_ {\mathbf {U}} f (z) d z
\]

is a continuous linear form on B(A).

\begin{enumerate}
\def\labelenumi{\alph{enumi})}
\setcounter{enumi}{1}
\tightlist
\item
  Let \(h\) be the finite-rank endomorphism of B(A) such that \(h(f) = \ell(f)\varphi_{\mathrm{A}}\) . The essential spectrum of \(u - h\) is distinct from that of \(u\) . (Show that it contains 0.)
\end{enumerate}

\begin{enumerate}
\def\labelenumi{\arabic{enumi})}
\setcounter{enumi}{2}
\tightlist
\item
  Let E be a complex Hilbert space. Let U be a connected open subset of \(\mathbf{C}\). Let f be a holomorphic function from U into the space \(\mathcal{L}^{c}(\mathrm{E})\) . Suppose that there exists \(z \in U\) such that 1 belongs to the resolvent set of \(f(z)\) .
\end{enumerate}

\begin{enumerate}
\def\labelenumi{\alph{enumi})}
\item
  The set D of \(z \in U\) such that \(1 \in \operatorname{Sp}(f(z))\) is a discrete subset of U.
\item
  The map from U-D to \(\mathcal{L}(\mathrm{E})\) defined by \(z\mapsto \mathrm{R}(f(z),1)\) is holomorphic.
\end{enumerate}

We fix a point \(z_0 \in \mathrm{D}\) .

\begin{enumerate}
\def\labelenumi{\alph{enumi})}
\setcounter{enumi}{2}
\item
  There exists an integer \(k \geqslant 1\) such that the map from U-D into \(\mathcal{L}(\mathrm{E})\) defined by \(z \mapsto (z - z_{0})^{k}\mathrm{R}(f(z), 1)\) extends to a holomorphic map on \((\mathrm{U}-\mathrm{D}) \cup \{z_{0}\}\) .
\item
  There exists a unique sequence \((u_{n})_{n\in\mathbf{N}}\) of endomorphisms of E such that \(u_{n}=0\) for n sufficiently large and a unique holomorphic function S from \((\mathrm{U}-\mathrm{D})\cup\{z_{0}\}\) into \(\mathcal{L}(\mathrm{E})\) such that for all \(z\in U-D\) , we have
\end{enumerate}

\[
\mathrm{R} (f (z), 1) = \sum_ {j = - k} ^ {- 1} (z - z _ {0}) ^ {j} u _ {j} + \mathrm{S} (z).
\]

\begin{enumerate}
\def\labelenumi{\alph{enumi})}
\setcounter{enumi}{4}
\tightlist
\item
  The endomorphism \(u_{-1}\) is the spectral projector of \(f(z_0)\) relative to the open and closed subset \(\{1\}\) of the spectrum of \(f(z_0)\) .
\end{enumerate}

\begin{enumerate}
\def\labelenumi{\arabic{enumi})}
\setcounter{enumi}{3}
\tightlist
\item
  Let E be a Banach space of countable type. Let \(u \in \mathcal{L}(\mathrm{E})\) . Assume that the sequence \((\| u^k \|)_{k \in \mathbf{N}}\) is bounded.
\end{enumerate}

For every \(\lambda\in\mathrm{Sp}(u)\cap\mathbf{U}\) , let \(x_{\lambda}\in\mathrm{E}\) be an eigenvector of \(u\) of norm 1 corresponding to \(\lambda\) .

\begin{enumerate}
\def\labelenumi{\alph{enumi})}
\tightlist
\item
  For all \(\lambda_1\) , \(\lambda_2\) in \(\mathrm{Sp}(u) \cap \mathbf{U}\) and every integer \(k \in \mathbf{N}\) , we have
\end{enumerate}

\[
\left| \lambda_ {1} ^ {k} - \lambda_ {2} ^ {k} \right| \leqslant \left(1 + \| u ^ {k} \|\right) \| x _ {\lambda_ {1}} - x _ {\lambda_ {2}} \|.
\]

\begin{enumerate}
\def\labelenumi{\alph{enumi})}
\setcounter{enumi}{1}
\item
  There exists a real number c \textgreater{} 0 such that \(\|x_{\lambda_{1}} - x_{\lambda_{2}}\| \geqslant c\) for all \(\lambda_{1} \neq \lambda_{2}\) in \(\operatorname{Sp}(u) \cap \mathbf{U}\) .
\item
  The set \(\operatorname{Sp}(u) \cap \mathbf{U}\) is countable (“Jamison’s theorem”).
\item
  Jamison's theorem does not hold in general if E is not of countable type.
\end{enumerate}

\begin{enumerate}
\def\labelenumi{\arabic{enumi})}
\setcounter{enumi}{4}
\tightlist
\item
  The goal of this exercise is to give another proof of Lomonosov's theorem (cor. 2 of III, p. 13) in the case where the space E is a Banach space.
\end{enumerate}

Let E be a complex Banach space of dimension at least 2. Let h be a nonzero compact endomorphism of E and let A be the commutant of h in \(\mathcal{L}(\mathrm{E})\) . a) One of the following conditions is satisfied:

\begin{enumerate}
\def\labelenumi{(\roman{enumi})}
\item
  There exists \(y \in E\) such that Ay is a nonzero subspace that is not dense in E;
\item
  For every closed bounded subset C of A with nonempty interior, there exists a finite subset I of A such that \(h(\mathrm{C})\) is contained in the union of the sets \(u^{-1}(\mathrm{C})\) for \(u \in I\) .
\end{enumerate}

\begin{enumerate}
\def\labelenumi{\alph{enumi})}
\setcounter{enumi}{1}
\tightlist
\item
  Assume that condition (ii) of the previous question is satisfied. Let C be a closed bounded subset of E with nonempty interior. For every integer \(m \geqslant 1\) there exists a finite sequence \((u_{1}, \ldots, u_{m})\) in A such that
\end{enumerate}

\[
u _ {1} \circ \dots \circ u _ {m} \circ h ^ {m} (x _ {0}) \in \mathrm{C}.
\]

(Play table tennis.)

\begin{enumerate}
\def\labelenumi{\alph{enumi})}
\setcounter{enumi}{2}
\item
  Under the same hypothesis, show that the spectral radius of h is not zero (in the contrary case, consider an element \(x_{0} \in E\) such that \(\|h(x_{0})\| > \|h\|\) and the closed ball C with center \(x_{0}\) and radius \(\|h\|\) ; noting that \(\|( \lambda h)^{n}\|\) tends to 0 for every \(\lambda \in R\) , show that the conclusion of the preceding question would imply that 0 belongs to the closure of C).
\item
  For every endomorphism \(u\) of E that commutes with a nonzero compact endomorphism of E, there exists a nonzero closed subspace \(F \subset E\) different from E such that \(u(F) \subset F\) .
\end{enumerate}

\begin{enumerate}
\def\labelenumi{\arabic{enumi})}
\setcounter{enumi}{5}
\tightlist
\item
  Let A be a unital Z-algebra. Assume that A is a free Z-module. Let B be a subset of A that is a basis of A as a Z-module. We say that the pair \((A, B)\) is a natural algebra if the structure constants of A with respect to B (A, III, p. 10) belong to N and if the coefficients of \(1_{A}\) in the basis B belong to N. Moreover, we say that (A, B) is unital if \(1_{A}\) is an element of B.
\end{enumerate}

Let (A, B) be a natural algebra. Let \(B_{0}\) be the set of \(b \in B\) such that the coefficient of b in \(1_{A}\) is nonzero. Denote by \(\tau_{A}\) the unique homomorphism of Z-modules from A to Z such that, for every b in B, one has \(\tau_{\mathrm{A}}(b) = 1\) if \(b \in B_{0}\) , and \(\tau_{\mathrm{A}}(b) = 0\) otherwise.

\begin{enumerate}
\def\labelenumi{\alph{enumi})}
\item
  One has \(\tau_{\mathrm{A}}(xy) = \tau_{\mathrm{A}}(yx)\) for all \((x,y)\in \mathrm{A}^2\) .
\item
  For all \(b \neq b'\) in \(B_{0}\) , we have \(b^2 = b\) and \(bb' = 0\) .
\end{enumerate}

\begin{enumerate}
\def\labelenumi{\arabic{enumi})}
\setcounter{enumi}{6}
\tightlist
\item
  Let (A, B) be a natural algebra.
\end{enumerate}

\begin{enumerate}
\def\labelenumi{\alph{enumi})}
\tightlist
\item
  There exists at most one involution \(\iota\) of \(\mathcal{B}\) such that the unique \(\mathbf{Z}\)-module homomorphism from A to A sending \(b\) to \(\iota(b)\) is an anti-automorphism of A satisfying the condition \(\tau_{\mathrm{A}}(bb') = 1\) if \(b' = \iota(b)\) and \(\tau_{\mathrm{A}}(bb') = 0\) otherwise.
\end{enumerate}

When such an involution \(\iota\) exists, we say that (A, \(\mathcal{B}\) ) is a natural basal algebra.

We now assume that (A, B) is a natural basal algebra.

\begin{enumerate}
\def\labelenumi{\alph{enumi})}
\setcounter{enumi}{1}
\tightlist
\item
  We then have \(\iota(b) = b\) for all \(b \in B_{0}\) as well as the formula
\end{enumerate}

\[
1 _ {\mathrm{A}} = \sum_ {b \in \mathcal {B} _ {0}} b.
\]

\begin{enumerate}
\def\labelenumi{\alph{enumi})}
\setcounter{enumi}{2}
\tightlist
\item
  The structure constants \(\gamma_{bb'}^{b''}\) of A satisfy \(\gamma_{bb'}^{\iota(b'')} = \tau_{\mathrm{A}}(bb'b'')\) for \((b, b', b'')\) in \(\mathcal{B}^{3}\).
\end{enumerate}

\begin{enumerate}
\def\labelenumi{\arabic{enumi})}
\setcounter{enumi}{7}
\tightlist
\item
  Let (A, B) be a natural basal algebra such that A has finite rank.
\end{enumerate}

\begin{enumerate}
\def\labelenumi{\alph{enumi})}
\tightlist
\item
  For every \(x \in A\) , the element
\end{enumerate}

\[
\sum_ {b \in \mathcal {B}} b x \iota (b)
\]

is central.

In the rest of this exercise, we further assume that (A, \(\mathcal{B}\) ) is unital.

\begin{enumerate}
\def\labelenumi{\alph{enumi})}
\setcounter{enumi}{1}
\item
  For all \(b\) and \(b'\) in \(\mathcal{B}\) , there exist \(b_1\) and \(b_2\) in \(\mathcal{B}\) such that the coefficient of \(b'\) in \(bb_1\) and in \(b_2b\) is nonzero.
\item
  For every \(x \in B\) , the matrix representing, in the basis B, the map \(y \mapsto xy\) from A to A is a real matrix all of whose coefficients are positive. We denote by pf the unique homomorphism from A to R such that \(\operatorname{pf}(b)\) is the spectral radius of this matrix when \(b \in B\) .
\end{enumerate}

\(d)\) For every \(b \in \mathcal{B}\) , the real number \(\alpha = \mathrm{pf}(b)\) is an algebraic integer such that \(\alpha \geqslant 1\) ; for every conjugate \(\alpha' \in \mathbf{C}\) of \(\alpha\) , we have \(|\alpha'| \leqslant \alpha\) .

\begin{enumerate}
\def\labelenumi{\alph{enumi})}
\setcounter{enumi}{4}
\tightlist
\item
  Let \(a_0 \in \mathrm{A}\) be the sum of the elements of \(\mathcal{B}\) . All coefficients of the matrix representing the map \(f: y \mapsto ya_0\) in the basis \(\mathcal{B}\) are strictly positive. There exists an element \(r \in \mathbf{C} \otimes \mathrm{A}\) all of whose coefficients in the basis \((1 \otimes b)_{b \in \mathcal{B}}\) are strictly positive such that \(f_{\mathbf{C}}(r) = \varrho(f_{\mathbf{C}})r\) .
\end{enumerate}

\(f)\) For every \(x \in \mathrm{A}\) , there exists a unique complex number \(\delta(x)\) such that \(xr = \delta(x)r\) in \(\mathbf{C} \otimes \mathrm{A}\) . The map \(\delta: \mathrm{A} \to \mathbf{C}\) is a character of \(\mathrm{A}\) .

\begin{enumerate}
\def\labelenumi{\alph{enumi})}
\setcounter{enumi}{6}
\tightlist
\item
  One has \(\delta(x) = \mathrm{pf}(x)\) for all \(x \in \mathrm{A}\) . In particular, the map \(x \mapsto \mathrm{pf}(x)\) is a character of A.
\end{enumerate}

\(h)\) The map pf is the unique character of A such that \(\mathrm{pf}(b) \geqslant 0\) for all \(b \in \mathcal{B}\) ; moreover, one has \(\mathrm{pf}(b) > 0\) for all \(b \in \mathcal{B}\) .

\begin{enumerate}
\def\labelenumi{\roman{enumi})}
\tightlist
\item
  For every \(y \in \mathrm{A}\) , we have \(ry = \delta(y)r\) in \(\mathbf{C} \otimes \mathbf{A}\) .
\end{enumerate}

\begin{enumerate}
\def\labelenumi{\alph{enumi})}
\setcounter{enumi}{9}
\tightlist
\item
  The element
\end{enumerate}

\[
r _ {0} = \sum_ {b \in \mathcal {B}} \mathrm{pf} (b) b
\]

possesses the properties required in question i).

\begin{enumerate}
\def\labelenumi{\arabic{enumi})}
\setcounter{enumi}{8}
\tightlist
\item
  \begin{enumerate}
  \def\labelenumii{\alph{enumii})}
  \tightlist
  \item
    Let \(n\) be an integer and \(M\) a square matrix with coefficients in \(\mathbf{N}\) . One assumes that \(\varrho(M^t M) = \varrho(M)^2\) and that \(|\varrho(M)| < 2\) . There exists an integer \(m \geqslant 2\) such that \(\varrho(M) = 2\cos(\pi/m)\) . (Use Kronecker's theorem, cf. A, V, p. 155, exercise 17.)
  \end{enumerate}
\end{enumerate}

\begin{enumerate}
\def\labelenumi{\alph{enumi})}
\setcounter{enumi}{1}
\item
  For every integer \(m \geqslant 2\) , there exists a square matrix M with coefficients in N such that \(\varrho(M) = 2 \cos(\pi/m)\) .
\item
  Let (A, B) be a unital natural basal algebra of finite rank. Let \(b \in B\) . If \(\operatorname{pf}(b) < 2\) , then there exists an integer \(m \geqslant 3\) such that \(\operatorname{pf}(b) = 2 \cos(\pi/m)\) .
\end{enumerate}

Exercises 10 to 24 provide another proof of Corollary 5 of III, p. 100, as well as certain refinements. The following notions will be used there.

Let I be a finite set and let n be its cardinality.

We denote by \(1_{\mathrm{I}}\) the real identity matrix of type (I, I).

- A square matrix of type (I, I) with real coefficients is said to be positive if all its coefficients are \(\geqslant 0\) and strictly positive if all its coefficients are \(>0\) .

- A square matrix \((a_{ij})_{i,j\in \mathrm{I}}\) of type (I,I) with real coefficients is said to be reducible if there exists a partition of I into two nonempty subsets I' and I'\,' such that \(a_{ij} = 0\) for all \((i,j)\in \mathrm{I}'\times \mathrm{I}''\) ; it is said to be irreducible if it is not reducible.

- One calls a coordinate subspace of \(\mathbf{R}^{\mathrm{I}}\) any vector subspace of \(\mathbf{R}^{\mathrm{I}}\) generated by a subset of the vectors of the canonical basis. There are \(2^{n}\) coordinate subspaces; for each integer \(k\in\{0,\ldots,n\}\) , those among them of dimension \(k\) number \(\binom{n}{k}\) .

A real matrix of type (I, I) is therefore irreducible if and only if the only coordinate subspaces it stabilizes are \(\{0\}\) and \(R^{I}\) .

For all \(x\) and \(y\) in \(\mathbf{R}^{\mathrm{I}}\) we will write \(x \preccurlyeq y\) (and \(y \succcurlyeq x\) ) if \(y - x\) is in the cone of elements of \(\mathbf{R}^{\mathrm{I}}\) with positive coordinates, that is, if \(x_i \leqslant y_i\) for all \(i \in \mathrm{I}\) . We will write \(x \prec y\) (and \(y \succ x\) ) if all coordinates of \(y - x\) are strictly positive.

\begin{enumerate}
\def\labelenumi{\arabic{enumi})}
\setcounter{enumi}{9}
\item
  Let \(A\) be a real matrix of type (I, I), positive and irreducible. For every integer \(n\) such that \(n \geqslant \operatorname{Card}(\mathrm{I})\) , the matrix \((1_{\mathrm{I}} + A)^{n-1}\) is strictly positive. (Let \(y \in \mathbf{R}_{+}^{n}\) and \(z = (1_{\mathrm{I}} + A)y\) ; if \(y \neq 0\), prove that the set of \(i \in \mathrm{I}\) such that \(z_{i} = 0\) is strictly contained in \(\{i \in \mathrm{I} \mid y_{i} = 0\}\) .)
\item
  Let \(A\) be a real matrix of type (I, I), positive and irreducible. We denote by \(\psi \in \mathbf{R}[\mathrm{X}]\) the minimal polynomial of \(A\) and we set \(m = \deg(\psi)\) .
\end{enumerate}

For every integer \(q\), set \(A^q = (a_{ij}^{(q)})_{i,j \in \mathrm{I}}\) .

For every \((i,j)\in\mathrm{I}^{2}\) , there exists an integer q such that \(1\leqslant q\leqslant m\) and \(a_{ij}^{(q)}>0\) . (Let r be the remainder of the Euclidean division of \((1+\mathrm{X})^{n-1}\) by \(\psi\) ; consider the matrix \(r(A)\) and observe that it is strictly positive.)

\begin{enumerate}
\def\labelenumi{\arabic{enumi})}
\setcounter{enumi}{11}
\tightlist
\item
  Let A be a real matrix of type (I, I), irreducible and positive. For every \(x \in R^{I}\) non-zero, we set
\end{enumerate}

\[
r_{x} = \inf_{\substack{i\in \mathrm{I}\\ x_{i}\neq 0}}\frac{(Ax)_{i}}{x_{i}},
\]

where \((Ax)_i = \sum_{k\in I}a_{ik}x_k\) is the coordinate of \(Ax\) with index \(i\) .

\begin{enumerate}
\def\labelenumi{\alph{enumi})}
\item
  If \(x \succcurlyeq 0\) and \(x \neq 0\) , then \(r_x\) is the largest real number \(\varrho\) such that \(\varrho x \preccurlyeq Ax\) .
\item
  Suppose that \(x \succcurlyeq 0\) and set \(y = (1_{\mathrm{I}} + A)x\) . We then have \(r_x \leqslant r_y\) .
\item
  The restriction to \((\mathbf{R}_{+}^{*})^{\mathrm{I}}\) of the function \(x \mapsto r_{x}\) is continuous.
\end{enumerate}

\(d)\) Deduce that there exists \(z \in \mathbf{R}^{\mathrm{I}}\) such that \(z \succcurlyeq 0\) and such that \(r_z = \sup_{x \succcurlyeq 0} r_x\) (Use compact sets \(\mathrm{M} = \{x \in \mathbf{R}_{+}^{\mathrm{I}} \mid \sum_{i \in \mathrm{I}} x_i = 1\}\) as well as \((1_{\mathrm{I}} + A)^{n-1}(\mathrm{M})\) , then exercise 10.)

We set \(r = r_{z}\). An element z of \(R^{I}\) such that \(z \succcurlyeq 0\) and \(r_{z} = r\) will be said to be extremal.

\begin{enumerate}
\def\labelenumi{\alph{enumi})}
\setcounter{enumi}{4}
\item
  We have r \textgreater{} 0.
\item
  If z is an extremal vector, then Az = rz. (If \(u = Az - rz \succcurlyeq 0\) were nonzero, then \((1_{\mathrm{I}} + A)^{n-1}u\) would be \(\succ 0\) based on Exercise 10 and, for \(x = (1_{\mathrm{I}} + A)^{n-1}z\) , one would have \(Ax \succ rx\) .)
\item
  Every extremal element is \(\succ 0\) .
\end{enumerate}

For every vector \(y\) in \(\mathbf{C}^{\mathrm{I}}\) , we denote by \(y^{+}\) the vector \((|y_{i}|)_{i\in \mathrm{I}}\) ; we therefore have \(y^{+}\in \mathbf{R}^{\mathrm{I}}\) and \(y^{+}\succcurlyeq 0\) .

\begin{enumerate}
\def\labelenumi{\alph{enumi})}
\setcounter{enumi}{7}
\item
  Let \(\alpha \in \mathbf{C}\) be an eigenvalue of \(A\) , viewed as a complex matrix, and let \(y\) in \(\mathbf{C}^{\mathrm{I}}\) be a corresponding eigenvector. We have \(|\alpha|y^{+} \preccurlyeq Ay^{+}\) and \(|\alpha| \leqslant r\) .
\item
  Deduce that if \(\alpha = r\) all coordinates of y are nonzero, and then that the eigenspace (in \(C^{I}\) ) of A associated with the eigenvalue r has dimension 1.
\end{enumerate}

We denote by \(\Delta(\mathrm{X})\) the determinant of \(\mathrm{X} \cdot 1_{\mathrm{I}} - A\) and by \(B(\mathrm{X})\) the transpose of the matrix of cofactors of \(\mathrm{X} \cdot 1_{\mathrm{I}} - A\) . We write \(B(\mathrm{X}) = (B_{ij}(\mathrm{X}))_{i,j \in \mathrm{I}}\) . We therefore have \(B(\mathrm{X})(\mathrm{X} \cdot 1_{\mathrm{I}} - A) = (\mathrm{X} \cdot 1_{\mathrm{I}} - A)B(\mathrm{X}) = \Delta(\mathrm{X}) \cdot 1_{\mathrm{E}}\) (A, III, p. 99, prop. 12).

\begin{enumerate}
\def\labelenumi{\alph{enumi})}
\setcounter{enumi}{9}
\tightlist
\item
  For every \(i \in I\) , the column with index i of \(B(r)\) is nonzero. (If this is not the case, construct an eigenvector of A associated with the eigenvalue r whose i-th coefficient is zero, in contradiction with part (f)).
\end{enumerate}

\(k\) ) All the coefficients of \(B(r)\) have the same sign. (First show, using the relation \((r1_{\mathrm{I}} - A)B(r) = 0\) that for all \(i\in \mathrm{I}\) , all the coefficients of the column with index \(i\) of \(B(r)\) are nonzero and of the same sign, and then reason with the transposes.)

\(l)\) One has \(\Delta'(r) > 0\) and \(B_{ij}(r) > 0\) for all \((i,j) \in \mathrm{I}^2\) (Use the formula \(\Delta'(X) = \sum_i B_{ii}(X)\) and observe that \(\Delta'\) does not change sign on the interval \([r, +\infty]\) .)

\begin{enumerate}
\def\labelenumi{\alph{enumi})}
\setcounter{enumi}{12}
\tightlist
\item
  Conclude that r is a simple root of the characteristic polynomial of A.
\end{enumerate}

\begin{enumerate}
\def\labelenumi{\arabic{enumi})}
\setcounter{enumi}{12}
\tightlist
\item
  Let A (resp. C) be a real (resp. complex) square matrix of type (I, I). Assume that A is positive and irreducible, and that \(C^{+} \preccurlyeq A\) , where, similarly to Exercise 12, we denote by \(C^{+}\) the matrix whose coefficients are the absolute values of the coefficients of C, and we write \(C^{+} \preccurlyeq A\) to indicate that the coefficients of \(A - C^{+}\) are positive.
\end{enumerate}

We resume the notation of the preceding exercise.

\begin{enumerate}
\def\labelenumi{\alph{enumi})}
\item
  Let \(\gamma \in \mathbf{C}\) be an eigenvalue of \(C\) and let \(y\) in \(\mathbf{C}^{\mathrm{I}}\) be an associated eigenvector. We have \(|\gamma|y^{+} \preccurlyeq C^{+}y^{+} \preccurlyeq Ay^{+}\) and \(|\gamma| \leqslant r_{y^{+}} \leqslant r\) .
\item
  We further assume that \(|\gamma| = r\) . Then \(y^{+}\) is an extremal vector for A (cf. question f) of Exercise 12), and we have the relations \(Ay^{+} = C^{+}y^{+} = ry^{+}\) and \(A = C^{+}\) .
\end{enumerate}

Let \(u\) (resp. \(u_j\) , for \(j \in I\) ) be the unique complex number such that \(\gamma = |\gamma|u\) (resp. \(y_j = |y_j|u_j\) , for \(j \in I\) ). Let us denote by \(D\) the square complex matrix of type (I, I) whose diagonal coefficients are \((u_j)_{j \in I}\) and whose other coefficients are zero, so that \(y = Dy^+\). Let \(F = u^{-1}D^{-1}CD\) .

\begin{enumerate}
\def\labelenumi{\alph{enumi})}
\setcounter{enumi}{2}
\item
  One has \(Fy^{+} = ry^{+}\) and \(F^{+}y^{+} = Fy^{+}\) (Observe that \(F^{+} = C^{+}\) .)
\item
  One has \(F = F^{+}\) and \(C = uDAD^{-1}\) .
\end{enumerate}

\begin{enumerate}
\def\labelenumi{\arabic{enumi})}
\setcounter{enumi}{13}
\tightlist
\item
  Let A be a real square matrix of type (I, I), positive and irreducible. We resume the notation of Exercise 12.
\end{enumerate}

Let \((\lambda_{h})_{h\in\mathrm{H}}\) be the family of complex eigenvalues of A which have modulus r, repeated with their multiplicities; we set \(v_{h}=\lambda_{h}/r\) for all \(h\in H\) .

Let \(h \in \mathrm{H}\) ; by Exercise 13 applied to \(C = A\) and \(\gamma = \lambda_h\) , there exists a matrix \(D_{(h)}\) such that \(A = v_h D_{(h)} AD_{(h)}^{-1}\) and \(D_{(h)}^+ = 1_{\mathrm{I}}\) .

\begin{enumerate}
\def\labelenumi{\alph{enumi})}
\item
  Let \(z \succ 0\) be an extremal vector of A. For \(h \in H\) , the vector \(y_{(h)} = D_{(h)}z\) is an eigenvector of A for the eigenvalue \(\lambda_{h}\) .
\item
  Show that for \(h \in H\) , the eigenvalue \(\lambda_{h}\) of A is simple. Deduce that the matrices \(D_{(h)}\) and the vectors \(y_{(h)}\) are uniquely determined up to a scalar factor.
\item
  Show that, for all h and k in \(H^{2}\) , we have \(A = u_{h}u_{k}^{-1}D_{(h)}D_{(k)}^{-1}AD_{(k)}D_{(h)}^{-1}\) and hence that the vector \(D_{(h)}D_{(k)}^{-1}z\) is an eigenvector of A associated with the eigenvalue \(u_{h}u_{k}^{-1}r\) .
\item
  Deduce that \(\{u_h\}_{h\in \mathrm{H}}\) is a subgroup of \(\mathbf{C}^*\) and that there exists a family \((\mu_h)_{h\in \mathrm{H}}\) of nonzero complex numbers such that \(\{\mu_h^{-1}D_{(h)}\}_{h\in \mathrm{H}}\) is a subgroup of the group of invertible matrices of type (I,I) with complex coefficients, and that these subgroups are cyclic of order Card(H). (Choose an index \(i\in \mathrm{I}\) and normalize \(\mathrm{D}_{(h)}\) so that its coefficient with row and column index \(i\) is equal to 1.)
\end{enumerate}

In the rest of the exercise, we assume that H = Z/dZ and that there exists \(\omega \in C^{*}\) such that \(\omega^{d} = 1\) and such that \(u_{h} = \omega^{h}\) for \(h \in H\) . Setting \(D = D_{(1)}\) , we then have \(D^{d} = 1_{I}\) and \(D_{(h)} = D^{h}\) for all \(h \in H\) .

\begin{enumerate}
\def\labelenumi{\alph{enumi})}
\setcounter{enumi}{4}
\tightlist
\item
  Deduce that \(A = \omega DAD^{-1}\) .
\end{enumerate}

We consider the partition \((\mathrm{I}_{h})_{h\in\mathrm{H}}\) of I, so that for \(i\in I\) , the coefficient with indices \((i,i)\) of D is equal to \(\omega^{h}\) if and only if \(i\in I_{h}\) . For every h, we set \(n_{h}=\mathrm{Card}(I_{h})\) . The matrix D is written “in blocks”

\[
D = \left( \begin{array}{c c c c} 1 _ {n _ {0}} & & & \\ & \omega 1 _ {n _ {1}} & & \\ & & \ddots & \\ & & & \omega^ {d - 1} 1 _ {n _ {d - 1}} \end{array} \right)
\]

and we consider the corresponding decomposition of \(A\) :

\[
A = \left( \begin{array}{c c c c} A _ {0, 0} & A _ {0, 1} & \dots & A _ {0, d - 1} \\ A _ {1, 0} & A _ {1, 1} & \dots & A _ {1, d - 1} \\ \vdots & \vdots & \ddots & \vdots \\ A _ {d - 1, 0} & A _ {d - 1, 1} & \dots & A _ {d - 1, d - 1} \end{array} \right),
\]

so that for \(h, k \in \mathrm{H}\) , \(A_{h,k}\) is a matrix of type \((\mathrm{I}_h, \mathrm{I}_k)\) .

\(f)\) For all \((h,k)\in\mathrm{H}\) , we have \(\omega A_{h,k}=\omega^{k-h}A_{h,k}\) . Deduce that \(\mathrm{A}_{h,k}=0\) if \(k-h\neq1\) .

\begin{enumerate}
\def\labelenumi{\arabic{enumi})}
\setcounter{enumi}{14}
\tightlist
\item
  We resume the hypotheses and notation of Exercise 12. Let \(\Delta_1(\mathrm{X})\) be the greatest common divisor of the coefficients of \(B(\mathrm{X})\) ; the polynomial \(\Delta_1(\mathrm{X})\) is assumed to be monic.
\end{enumerate}

\begin{enumerate}
\def\labelenumi{\alph{enumi})}
\item
  The polynomial \(\Delta_1(X)\) divides \(\Delta(X)\) and \(\Delta_1(r) \neq 0\) .
\item
  Deduce that the real roots of \(\Delta_{1}(X)\) are strictly smaller than r. Conclude that \(\Delta_{1}(r)>0\) .
\item
  Set \(C(\mathrm{X}) = \frac{1}{\Delta_1(\mathrm{X})} B(\mathrm{X})\) . This is a matrix with coefficients in \(\mathbf{R}[\mathrm{X}]\) . Show that \(C(r) \succ 0\) .
\end{enumerate}

\begin{enumerate}
\def\labelenumi{\arabic{enumi})}
\setcounter{enumi}{15}
\tightlist
\item
  Let A be a real square matrix of type (I, I), positive and irreducible. For all \(j \in I\) set \(s_{j} = \sum_{k \in I} a_{jk}\) ,
\end{enumerate}

\[
s = \inf _ {j \in \mathrm{I}} s _ {j}, \quad \text {and} \quad \mathrm{S} = \sup _ {j \in \mathrm{I}} s _ {j}.
\]

We also adopt the notation \(r\) and \(B(\mathbf{X})\) of Exercise 12.

\begin{enumerate}
\def\labelenumi{\alph{enumi})}
\item
  One has \(\sum_{j\in \mathrm{I}}(r - s_j)B_{kj}(r) = 0\)
\item
  Deduce that \(s \leqslant r \leqslant S\) and that the equality case in either of these two inequalities implies that s = S.
\end{enumerate}

\begin{enumerate}
\def\labelenumi{\arabic{enumi})}
\setcounter{enumi}{16}
\item
  Let \(A\) be a real square matrix of type (I, I), positive and irreducible. Let \(y \in \mathbf{R}^{\mathrm{I}}\) be an eigenvector of \(A\) associated with an eigenvalue \(\alpha\) such that \(y \succcurlyeq 0\) . We then have \(\alpha = r\) (with the notations of Exercise 12) and \(y \succ 0\) . (Consider an extremal vector \(u\) for \(^{t}A\) and compute \(\langle y|^{t}Au\rangle = \langle Ay|u\rangle\) .)
\item
  Let A be a real square matrix of type (I, I), positive and irreducible. Let \(\varrho\) be the spectral radius of A. For every \(x \in R^{I}\) such that \(y \succcurlyeq 0\) and \(x \neq 0\) we set
\end{enumerate}

\[
\varrho^ {x} = \sup _ {i \in \mathrm{I}} \frac {(A x) _ {i}}{x _ {i}},
\]

(with the convention \(\varrho^x = +\infty\) if there exists an index \(i\) such that \(x_i = 0\) and \((Ax)_i \neq 0\) ). Let

\[
\varrho^{\prime} = \inf_{\substack{x\succ 0\\ x\neq 0}}\varrho^{x}.
\]

Following an approach similar to that of Exercise 12, show that \(\varrho'\) is an eigenvalue of \(A\) and that \(\varrho' = \varrho\) .

\begin{enumerate}
\def\labelenumi{\arabic{enumi})}
\setcounter{enumi}{18}
\tightlist
\item
  Let A be a real square matrix of type (I, I), positive and irreducible; let \(\varrho\) be its spectral radius. Let \(z \in R^{I}\) such that \(z \succcurlyeq 0\) and \(z \neq 0\) .
\end{enumerate}

If \(\varrho z \preccurlyeq Az\) (resp. \(\varrho z \succcurlyeq Az\) ), then \(\varrho z = Az\) and \(z \succ 0\) .

\begin{enumerate}
\def\labelenumi{\arabic{enumi})}
\setcounter{enumi}{19}
\tightlist
\item
  Let A be a real square matrix of type (I, I), positive but not necessarily irreducible; let \(\varrho\) be the spectral radius of A. We keep the notations \(B(\mathrm{X})\) and \(C(\mathrm{X})\) from the preceding exercises.
\end{enumerate}

\begin{enumerate}
\def\labelenumi{\alph{enumi})}
\item
  The real number \(\varrho\) is an eigenvalue of \(A\) and there exists at least one eigenvector \(\succcurlyeq 0\) associated with \(\varrho\) .
\item
  For every real number \(\lambda \geqslant r\) , the matrices \(B(\lambda)\) and \(B'(\lambda)\) are positive. (For the second matrix, one may proceed by induction on the cardinality of I.)
\item
  The matrix \(C(\lambda)\) is positive for all \(\lambda \geqslant r\) .
\item
  If A is irreducible, then for every real number \(\lambda > r\) , the matrices \(B(\lambda)\) , \(C(\lambda)\) and \((\lambda1_{\mathrm{I}} - A)^{-1}\) are strictly positive.
\item
  Let J be a subset of I different from I. Let \(\varrho_{\mathrm{J}}\) be the spectral radius of the submatrix \((a_{i,j})_{i,j\in \mathrm{J}}\) . We have \(\varrho_{\mathrm{J}}\leqslant \varrho\) and this inequality is strict if \(A\) is irreducible. Conversely, if \(A\) is reducible, there exists a subset \(\mathrm{J}\subset \mathrm{I}\) different from I, such that \(\varrho_{\mathrm{J}} = \varrho\) .
\item
  Let J be a subset of I different from I and such that the spectral radius of the submatrix \((a_{i,j})_{i,j\in\mathrm{J}}\) is equal to \(\varrho\) . For every subset K such that \(J\subset K\subset I\) , the spectral radius of the matrix \((a_{i,j})_{i,j\in\mathrm{K}}\) is equal to \(\varrho\) . Deduce that \(A\succ0\) if and only if all the diagonal coefficients of \(B(r)\) are strictly positive.
\end{enumerate}

\(g)\) For every real number \(\lambda >\varrho\) and for every subset J of I such that \(\mathrm{J}\neq \mathrm{I}\) the determinant of the matrix \((\lambda \delta_{i,j} - a_{i,j})_{i,j\in \mathrm{J}}\) is strictly positive.

\begin{enumerate}
\def\labelenumi{\alph{enumi})}
\setcounter{enumi}{7}
\tightlist
\item
  Let \(\lambda\) be a real number; we assume that the determinant of \((\lambda \delta_{i,j} - a_{i,j})_{i,j \in \mathbf{J}}\) is strictly positive for every subset \(\mathbf{J} \subset \mathbf{I}\) different from \(\mathbf{I}\) . We then have \(\lambda > \varrho\) .
\end{enumerate}

In the following questions, we assume that \(\mathrm{I}\) is the interval \([1,n]\) in \(\mathbf{N}\). i) Let \((g_{i,j})_{1\leqslant i,j\leqslant n}\) be a real matrix such that \(g_{i,j}\leqslant0\) for all \(i\neq j\) and such that the determinants of the matrices \((g_{i,j})_{1\leqslant i,j\leqslant m}\) are strictly positive for every integer \(m\) such that \(1\leqslant m\leqslant n\) . Then we have \(\det(g_{i,j})_{i,j\in\mathrm{J}}>0\) for every subset J of I.

\begin{enumerate}
\def\labelenumi{\alph{enumi})}
\setcounter{enumi}{9}
\item
  Let \(\lambda\) be a real number; we suppose that for every integer \(m\) such that \(1 \leqslant m \leqslant n\) , the determinant \(\det(\lambda \delta_{i,j} - a_{i,j})_{1 \leqslant i,j \leqslant m}\) is strictly positive. Then \(\lambda > \varrho\) .
\item
  Let \(C = (c_{i,j})_{1 \leqslant i,j \leqslant n}\) be a real matrix such that \(c_{i,j} \geqslant 0\) for all \(i \neq j\) . If \((-1)^{k} \det(c_{i,j})_{1 \leqslant i,j \leqslant m}\) is strictly positive for every integer m such that \(1 \leqslant m \leqslant n\) , then the real parts of the eigenvalues of C (viewed as a complex matrix) are negative.
\end{enumerate}

\begin{enumerate}
\def\labelenumi{\arabic{enumi})}
\setcounter{enumi}{20}
\tightlist
\item
  Let A be a real square matrix of type (I, I), positive; let \(\varrho\) be its spectral radius.
\end{enumerate}

\begin{enumerate}
\def\labelenumi{\alph{enumi})}
\tightlist
\item
  Prove that there exist integers p and q such that \(0 \leqslant q \leqslant p\) and a partition \((\mathrm{I}_{s})_{1 \leqslant s \leqslant p}\) of I into nonempty parts such that the corresponding block decomposition \(A = (A_{s,t})_{1 \leqslant s, t \leqslant p}\) has the following property:
\end{enumerate}

- The matrices \(A_{s,s}\) are irreducible;

- We have \(A_{s,t} = 0\) if \(t > s\) ;

- We have \(A_{s,t} = 0\) if \(1 \leqslant t < s \leqslant q\) ;

- For every integer \(s\) such that \(q \leqslant s \leqslant p\) , there exists \(t\) such that \(A_{s,t} \neq 0\) .

\begin{enumerate}
\def\labelenumi{\alph{enumi})}
\setcounter{enumi}{1}
\item
  Determine all possibilities for such a partition. (In particular, prove that the integers q and p are uniquely determined, as are the sets \(\{I_{1},\ldots,I_{g}\}\) and the sets \(\{I_{q+1},\ldots,I_{p}\}\) .)
\item
  There exists an eigenvector \(z \succ 0\) of A associated with the eigenvalue \(\varrho\) if and only if, for every integer \(s \in [1, p]\) the spectral radius of \(A_{s,s}\) is equal to \(\varrho\) if \(s \leqslant q\) and is strictly less than \(\varrho\) otherwise.
\item
  It is necessary and sufficient for both \(A\) and \({}^t A\) to admit a strictly positive extremal vector that there exists a partition as above such that \(A_{s,t} = 0\) for \(1 \leqslant t < s\) and \(q < s \leqslant p\) and such that the matrices \(A_{s,s}\) for \(s \leqslant q\) have spectral radius equal to \(\varrho\) .
\item
  If A and \(^{t}A\) both admit a strictly positive extremal vector, and if \(\varrho\) has spectral multiplicity equal to 1, then A is irreducible.
\end{enumerate}

\begin{enumerate}
\def\labelenumi{\arabic{enumi})}
\setcounter{enumi}{21}
\tightlist
\item
  Let A be a real square matrix of type (I, I), positive and irreducible; let \(\varrho\) be its spectral radius. Let \(d \geqslant 1\) be the integer determined by Exercise 14. We shall say that A is primitive if d = 1. Let \(\Delta(\mathrm{X})\) be the polynomial \(\det(\mathrm{X} \cdot 1_{\mathrm{I}} - A)\) .
\end{enumerate}

\begin{enumerate}
\def\labelenumi{\alph{enumi})}
\item
  We write \(\Delta(\mathrm{X}) = \mathrm{X}^{n} + a_{1}\mathrm{X}^{n_{1}} + \cdots + a_{t}\mathrm{X}^{n_{t}}\) , where \(a_{1}, \ldots, a_{t}\) are nonzero real numbers and \(n_{1}, \ldots, n_{t}\) natural integers such that \(n > n_{1} > \cdots > n_{t}\) . Prove that d is the greatest common divisor of \(n - n_{1}, n_{1} - n_{2}, \ldots, n_{t-1} - n_{t}\) .
\item
  There exists a nonzero integer m and a polynomial g such that \(\Delta(\mathrm{X}) = g(\mathrm{X}^{d})\mathrm{X}^{m}\) .
\item
  If there exists an integer p \textgreater{} 0 such that \(A^{p}\) is strictly positive, then A is primitive.
\item
  Suppose that \(A\) is primitive. Let us introduce the matrices \(C(\mathbf{X})\) of Exercise 15, and let us set, using the notation of that exercise, \(\mathsf{X}(\mathbf{X}) = \Delta (\mathbf{X}) / \Delta_1(\mathbf{X})\) . The limit of the sequence \((A^k /\varrho^k)_{k\in \mathbf{N}}\) is equal to \(C(\varrho) / \mathsf{X}'(\varrho)\) (One may use a decomposition of \(A\) into Jordan blocks.) Using the results of Exercise 15, deduce that there exists an integer \(p > 0\) such that \(A^p\) is strictly positive.
\item
  Suppose there exists an integer q \textgreater{} 0 such that \(A^{q}\) is reducible. Let p be the greatest common divisor of q and d. Prove that there exists a partition \((\mathrm{I}_{s})_{1 \leqslant s \leqslant p}\) of I into nonempty parts for which \(A^{q}\) is block diagonal, with irreducible diagonal blocks and spectral radius \(\varrho^{q}\) .
\item
  If A is primitive, then \(A^{p}\) is irreducible for every p \textgreater{} 0.
\end{enumerate}

\(g)\) There exists a partition \((\mathrm{I}_s)_{1 \leqslant s \leqslant d}\) of I into nonempty parts for which \(A^d\) is block diagonal, with primitive diagonal blocks and spectral radius equal to \(\varrho^d\) .

\begin{enumerate}
\def\labelenumi{\arabic{enumi})}
\setcounter{enumi}{22}
\tightlist
\item
  Let P be a real square matrix of type (I, I); we say that P is stochastic if it is nonnegative and if the vector \((1, \ldots, 1) \in \mathbf{R}^{\mathrm{I}}\) is an eigenvector of P associated with the eigenvalue 1.
\end{enumerate}

Let A be a real square matrix of type (I, I). Prove that A is positive and admits a strictly positive extremal vector if and only if there exists a stochastic matrix P, a real number r \textgreater{} 0 and a diagonal matrix Z with strictly positive diagonal entries such that \(A = rZPZ^{-1}\) . In this case the spectral radius of \(A\) is equal to \(r\) and an extremal vector of \(A\) is \(Z(1, \ldots, 1)\) .

\begin{enumerate}
\def\labelenumi{\arabic{enumi})}
\setcounter{enumi}{23}
\tightlist
\item
  \begin{enumerate}
  \def\labelenumii{\alph{enumii})}
  \tightlist
  \item
    Let P be a real square stochastic matrix of order n. The real number 1 is a simple root of the minimal polynomial of P. (One may use Exercise 20.) Moreover, every root of modulus 1 of the minimal polynomial of P is simple.
  \end{enumerate}
\end{enumerate}

\begin{enumerate}
\def\labelenumi{\alph{enumi})}
\setcounter{enumi}{1}
\tightlist
\item
  Let A be a positive real matrix admitting an extremal vector \(z \succ 0\) and of spectral radius \(\varrho\) . All roots of modulus \(\varrho\) of the minimal polynomial of A are simple.
\end{enumerate}

In Exercises 25 to 30, we consider a real Fréchet space E and a closed convex cone C in E, assumed to have nonempty interior. The order relation on E induced by C is denoted \(\preccurlyeq_{\mathrm{C}}\) . Thus, for \(x\) , \(y\) in E, we have \(x \succcurlyeq_{\mathrm{C}} y\) if and only if \(x - y \in \mathrm{C}\) . We will write \(x \succ_{\mathrm{C}} y\) for the relation \(x - y \in \mathring{\mathrm{C}}\) .

We consider a continuous map u: E → E satisfying the following conditions:

\begin{enumerate}
\def\labelenumi{(\roman{enumi})}
\item
  The map u is positively homogeneous: we have \(u(\lambda x) = \lambda u(x)\) for every real number \(\lambda \geqslant 0\) ;
\item
  The map u is increasing with respect to C, that is, \(x \succsim_{C} y\) implies \(u(x) \succsim_{C} u(y)\) ;
\item
  For every \(x \in \mathrm{C} - \{0\}\) , we have \(u(x) \in \mathring{\mathrm{C}}\) .
\end{enumerate}

The goal of these exercises, taken from the course of P.-L. Lions at the Collège de France, 2020/2021, is to prove, under certain conditions, that there exists an element \(x \in \mathrm{C} - \{0\}\) and \(\lambda > 0\) such that \(x = \lambda u(x)\) ; this provides in particular another proof of the Krein-Rutman theorem (III, p. 94, th. 4 and III, p. 97, prop. 8).

For \(f \in \mathring{\mathrm{C}}\) and \(\lambda \in \mathbf{R}_{+}\) , one calls an upper solution (resp. lower solution) of the equation \(y = \lambda u(y) + f\) an element \(y \in \mathrm{C}\) such that \(y \succcurlyeq_{\mathrm{C}} \lambda u(y) + f\) (resp. such that \(y \preccurlyeq_{\mathrm{C}} \lambda u(y) + f\) ). We denote by \(U_{f}\) the set of \(\lambda \in \mathbf{R}_{+}\) such that the set of upper solutions of the equation \(y = \lambda u(y) + f\) is nonempty.

\begin{enumerate}
\def\labelenumi{\arabic{enumi})}
\setcounter{enumi}{24}
\tightlist
\item
  \begin{enumerate}
  \def\labelenumii{\alph{enumii})}
  \tightlist
  \item
    Prove that for every \(x \in C\) and every \(y \in \mathring{C}\) , there exists a real number \(\lambda \geqslant 0\) such that \(x \preccurlyeq_{C} \lambda y\) .
  \end{enumerate}
\end{enumerate}

\begin{enumerate}
\def\labelenumi{\alph{enumi})}
\setcounter{enumi}{1}
\tightlist
\item
  The set \(U_{f}\) is an interval that contains 0.
\end{enumerate}

We introduce the following conditions on the map u.

\begin{enumerate}
\def\labelenumi{(\roman{enumi})}
\setcounter{enumi}{3}
\item
  For every \(\lambda \in \mathbf{R}_{+}\) , every \(f \in \mathring{\mathrm{C}}\) , every sequence \((x_{n})\) of elements of C that satisfies \(x_{n+1} = \lambda u(x_{n}) + f\) for all \(n\) and that is either increasing and bounded above, or decreasing, is convergent;
\item
  The image under u of every bounded subset of E is relatively compact.
\item
  For every map \(\lambda \mapsto \varepsilon_{\lambda}\) from \(U_f\) into C such that \(\lim_{\lambda \to \lambda_m} \varepsilon_{\lambda} = 0\) , the set
\end{enumerate}

\[
\{z \in \mathrm{C} \mid \exists \lambda \in \mathrm{U}_{f}\text{ such that } z = \lambda u(z) + \varepsilon_{\lambda} \}
\]

is relatively compact in E.

\begin{enumerate}
\def\labelenumi{\arabic{enumi})}
\setcounter{enumi}{25}
\tightlist
\item
  We resume the preceding notations and hypotheses and we suppose in addition that condition (iv) is satisfied.
\end{enumerate}

\begin{enumerate}
\def\labelenumi{\alph{enumi})}
\item
  Let \(\lambda \in \mathrm{U}_f\) . Let us define a sequence \((x_n)_{n \in \mathbf{N}}\) by \(x_0 = 0\) and \(x_{n+1} = \lambda u(x_n) + f\) for all \(n \in \mathbf{N}\) . It is increasing and bounded above; we have \(x_n \in \mathring{\mathrm{C}}\) for \(n \geqslant 1\) .
\item
  By condition (iv), the sequence \((x_{n})\) converges; let \(x(\lambda)\) denote its limit. Prove that we have \(x(\lambda) \in \mathring{\mathrm{C}}\) and \(x(\lambda) = \lambda u(x(\lambda)) + f\) .
\item
  Let \(\lambda \in \mathrm{U}_f\) and let \(y \in \mathrm{C}\) such that \(y \succcurlyeq_{\mathrm{C}} \lambda u(y) + f\) ; the sequence defined by \(x_0 = y\) , \(x_{n+1} = \lambda u(x_n) + f\) is decreasing, with values in C, and its limit \(x\) satisfies the equation \(x = \lambda u(x) + f\) .
\end{enumerate}

\(d)\) The upper bound of \(\mathrm{U}_{f}\) in \(\mathbf{R}_{+}\) does not belong to \(\mathrm{U}_{f}\) .

\begin{enumerate}
\def\labelenumi{\alph{enumi})}
\setcounter{enumi}{4}
\item
  One has \(U_{g} = U_{f}\) for all \(g \in \mathring{C}\) (For every \(\lambda \in U_{f}\) , let \(x \in C\) be such that \(x \succsim_{C} \lambda u(x) + f\) ; then there exists M \textgreater{} 0 such that y = Mx satisfies \(y \succsim_{C} \lambda u(y) + g\) .) f) Let \(\lambda \leqslant \mu\) be real numbers. Let \(x\) be a lower solution of the equation \(y = \lambda u(y) + f\) and let \(x'\) be an upper solution of the equation \(y = \mu u(y) + f\) . We then have \(x \preccurlyeq_{C} x'\) . (Show that the set of \(\varepsilon \in [0,1]\) such that \(\varepsilon x \preccurlyeq_{C} x'\) coincides with {[}0,1{]}.)
\item
  If C is salient, then for every \(\lambda \in \mathrm{U}_f\) , there exists a unique element \(x \in \mathring{\mathrm{C}}\) such that \(x = \lambda u(x) + f\) .
\end{enumerate}

\(h\) ) Let \(\mu\in\mathbf{R}_{+}\) such that \(f\preccurlyeq_{\mathrm{C}}\mu u(f)\) . We then have \(\mu\notin\mathrm{U}_{f}\) . (The sequence \((x_{n})_{n\in\mathbf{N}}\) defined by \(x_{0}=0\) and \(x_{n+1}=\mu u(x_{n})+f\) satisfies \(x_{n+1}-x_{n}\succcurlyeq_{\mathrm{C}}f\) for all \(n\in\mathbf{N}\) .) In particular, the interval \(\mathrm{U}_{f}\) is bounded. Henceforth we denote by \(\lambda_{\mathrm{m}}\) the supremum of \(\mathrm{U}_{f}\) .

\begin{enumerate}
\def\labelenumi{\roman{enumi})}
\tightlist
\item
  Let \(z \in \mathrm{C} - \{0\}\) and \(\mu \in \mathbf{R}_+\) be such that \(\mu u(z) \succsim_{\mathrm{C}} z\) . We then have \(\lambda_{\mathrm{m}} \leqslant \mu\) , and hence
\end{enumerate}

\[
\lambda_ {\mathrm{m}} \leqslant \inf _ {z \in \mathrm{C} - \{0 \}} \inf \{\mu \mid z \leqslant \mu u (z) \}.
\]

\begin{enumerate}
\def\labelenumi{\alph{enumi})}
\setcounter{enumi}{9}
\tightlist
\item
  The map \(\lambda \mapsto x(\lambda)\) (where \(x(\lambda)\) is defined in the question \(a\) ) is increasing.
\end{enumerate}

\(k)\) The map \(\lambda \mapsto x(\lambda)\) has no limit when \(\lambda\) tends to \(\lambda_{\mathrm{m}}\) . More precisely, if \((\lambda_n)_n\) is a sequence of \(\mathrm{U}_f\) which tends to \(\lambda_{\mathrm{m}}\) , then the sequence \((x(\lambda_n))_n\) is not convergent.

\begin{enumerate}
\def\labelenumi{\arabic{enumi})}
\setcounter{enumi}{26}
\tightlist
\item
  We assume that hypotheses (i) through (v) are satisfied.
\end{enumerate}

\begin{enumerate}
\def\labelenumi{\alph{enumi})}
\item
  Prove that the set of \(x(\lambda)\) , for \(\lambda \in \mathrm{U}_f\) , is not bounded.
\item
  Let d be a metric on E that defines its topology. For \(\lambda \in R_{+}\) , we define \(\hat{x}(\lambda) = d(0, x(\lambda))^{-1}x(\lambda)\) . Prove that the map \(\hat{x}\) is bounded and that the set of values of \(\hat{x}\) has a cluster value \(x_{m}\) such that
\end{enumerate}

\begin{enumerate}
\def\labelenumi{(\roman{enumi})}
\item
  \(d(0, x_{\mathrm{m}}) = 1,\)
\item
  \(x_{\mathrm{m}} = \lambda_{\mathrm{m}} u(x_{\mathrm{m}}),\)
\item
  \(x_{m} \in \overset{\circ}{C}.\)
\end{enumerate}

\begin{enumerate}
\def\labelenumi{\arabic{enumi})}
\setcounter{enumi}{27}
\tightlist
\item
  We assume in this exercise that conditions (i) through (iv) are satisfied. a) We fix \(x_0 \in \mathring{\mathrm{C}}\) . For \(\lambda \in \mathrm{U}_f\) , let \(\mathrm{A}(\lambda) = \inf \{\mathrm{M} \in \mathbf{R}_+ \mid x(\lambda) \preccurlyeq_{\mathrm{C}} \mathrm{M}x_0\}\) . Let \(\mathrm{C}_0 \in \mathbf{R}_+\) be such that \(f \preccurlyeq_{\mathrm{C}} \mathrm{C}_0 u(x_0)\) . For every \(\lambda\) in \(\mathrm{U}_f\) , the interval \([\lambda, \lambda + \mathrm{C}_0 / \mathrm{A}(\lambda)[\) is contained in \(\mathrm{U}_f\) (cf. exercise 26, question \(d\) )).
\end{enumerate}

\begin{enumerate}
\def\labelenumi{\alph{enumi})}
\setcounter{enumi}{1}
\tightlist
\item
  Hence deduce that the function \(\lambda \mapsto \mathrm{A}(\lambda)\) tends to \(+\infty\) as \(\lambda \in \mathrm{U}_f\) tends to \(\lambda_{\mathrm{m}}\) .
\end{enumerate}

For every \(\lambda\) in \(\mathrm{U}_f\) , we define \(\hat{x} (\lambda) = x(\lambda) / \mathrm{A}(\lambda)\) .

We now assume that condition (vi) is satisfied.

\begin{enumerate}
\def\labelenumi{\alph{enumi})}
\setcounter{enumi}{2}
\tightlist
\item
  Prove that the set of values of the mapping \(\hat{x}\) has a cluster value \(x_{\mathrm{m}}\) such that
\end{enumerate}

\[
0 \preccurlyeq_ {\mathrm{C}} x _ {\mathrm{m}} \preccurlyeq_ {\mathrm{C}} x _ {0} \quad \text {and} \quad x _ {\mathrm{m}} = \lambda_ {\mathrm{m}} u (x _ {\mathrm{m}}).
\]

\begin{enumerate}
\def\labelenumi{\alph{enumi})}
\setcounter{enumi}{3}
\item
  Prove that \(x_{m} \neq 0\) . (If one had \(x_{m} = 0\) , let U be a neighborhood of 0 in E such that \(x_{0} + U \subset C\) ; for every \(\mu \geqslant 0\) , there would exist a sequence \((\lambda_{n})\) in \(U_{f}\) converging to \(\lambda_{m}\) such that \(x_{0} - \mu \hat{x}(\lambda_{n}) \succsim_{C} 0\) , contradicting the definition of \(A(\lambda)\) .)
\item
  Let \(y \in \mathrm{C} - \{0\}\) be such that \(y = \lambda_{\mathrm{m}} u(y)\) . Set \(\theta = \sup \{\gamma \mid \gamma y \preccurlyeq_{\mathrm{C}} x_{\mathrm{m}}\}\) . We have \(\theta \in \mathbf{R}_{+}^{*}\) and \(x_{\mathrm{m}} = \theta y\) . (We have \(\theta y \preccurlyeq_{\mathrm{C}} x_{\mathrm{m}}\) and \(\theta y = \lambda_{\mathrm{m}} u(\theta y)\) ; if \(x_{\mathrm{m}} \neq \theta y\) , then we would have \(x_{\mathrm{m}} - \theta y \succ_{\mathrm{C}} 0\) and there would exist \(\varepsilon > 0\) such that \(x_{\mathrm{m}} - \theta y \succcurlyeq_{\mathrm{C}} \varepsilon y\) , contradicting the definition of \(\theta\) .)
\item
  Let \(y \in C - \{0\}\) and \(\mu \in R\) be such that \(y = \mu u(y)\) . We then have \(\mu = \lambda_{m}\) . (Note that \(0 < \mu\) and \(\mu \geqslant \lambda_{m}\) ; if we had \(\mu > \lambda_{m}\) , let \(\theta = \sup\{\gamma \mid \gamma x_{m} \preccurlyeq_{C} y\}\) ; we would then have \(\theta \in R_{+}^{*}\) , then \(\theta x_{m} \preccurlyeq_{C} y\) , whence \(\theta x_{m} = \lambda_{m} u(\theta x_{m}) \prec_{C} \mu u(\theta x_{m}) \preccurlyeq_{C} \mu u(y) = y\) , contradicting the definition of \(\theta\) .)
\end{enumerate}

\begin{enumerate}
\def\labelenumi{\arabic{enumi})}
\setcounter{enumi}{28}
\tightlist
\item
  We resume the notation of the preceding exercise; in particular, we assume that C has nonempty interior. We assume that u is a linear map and that conditions (i), (ii), and (iii) are satisfied.
\end{enumerate}

\begin{enumerate}
\def\labelenumi{\alph{enumi})}
\tightlist
\item
  If u is a compact endomorphism of E, then conditions (iv) and (v) are satisfied.
\end{enumerate}

We assume that conditions (iv) and (vi) are satisfied. Let \((\lambda_{\mathrm{m}}, x_{\mathrm{m}})\) be given by the preceding exercise.

\begin{enumerate}
\def\labelenumi{\alph{enumi})}
\setcounter{enumi}{1}
\tightlist
\item
  For every \(\lambda \in \mathbf{R}_{+}\) such that \(\lambda < \lambda_{\mathrm{m}}\) , the endomorphism \(1_{\mathrm{E}} - \lambda u\) is invertible. (First show that \(1_{\mathrm{E}} - \lambda u\), induced by passage to subspaces, is a bijection of \(\mathring{\mathrm{C}}\) , cf. questions \(a\) ) and \(f\) ) of the preceding exercise; then deduce that
\end{enumerate}

\(1_{E} - \lambda u\) is surjective, and then that it is injective by writing \(x_{1} - x_{2}\) as an element of the kernel of \(1_{E} - \lambda u\) , with \(x_{i} \in \mathring{C}\) , and verifying that there exists \(z \in E\) such that \(x_{1} + z - \lambda u(x_{1} + z) \in \mathring{C}\) .

\begin{enumerate}
\def\labelenumi{\alph{enumi})}
\setcounter{enumi}{2}
\tightlist
\item
  Let y in E be such that \(y = \lambda_{\mathrm{m}} u(y)\) . We then have \(y \in C\) . (There exists c such that \(cx_{m} + y \succ_{C} 0\) ; conclude using part (e) of the preceding exercise.)
\end{enumerate}

\(d)\) Let \(\mathrm{C}^{\circ}\subset\mathrm{E}^{\prime}\) be the polar of C. For every \(x\in\mathring{\mathrm{C}}\) and for every \(\xi\in\mathrm{C}^{\circ}-\{0\}\) , we have \(\langle x,\xi\rangle>0\) .

\begin{enumerate}
\def\labelenumi{\alph{enumi})}
\setcounter{enumi}{4}
\tightlist
\item
  The polar C° of C is preserved by \(^{t}u\colon E'\to E'\) ; for every \(\lambda\in R_{+}\) such that \(\lambda<\lambda_{m}\) , the endomorphism \(1_{E'}-\lambda^{t}u\) of \(E'\) is an isomorphism.
\end{enumerate}

\begin{enumerate}
\def\labelenumi{\arabic{enumi})}
\setcounter{enumi}{29}
\tightlist
\item
  We resume the hypotheses of the preceding exercise, so that u is linear, compact, and conditions (i) to (v) are assumed to be satisfied. The endomorphism \(^{t}u\) is also compact.
\end{enumerate}

\begin{enumerate}
\def\labelenumi{\alph{enumi})}
\item
  Let \(\eta \in \mathrm{C}^{\circ} - \{0\}\) . For every \(\lambda \in \mathbf{R}_{+}\) such that \(\lambda < \lambda_{\mathrm{m}}\) , the unique solution \(\xi(\lambda)\) of the equation \(\xi - \lambda^{t}u(\xi) = \eta\) belongs to \(\mathrm{C}^{\circ}\) . (If \(x \in \mathrm{C}\) , let \(y \in \mathrm{C}\) be such that \(y - \lambda u(y) = x\) ; we have \(\langle y, \eta \rangle = \langle x, \xi \rangle\) .)
\item
  Let \(x_{0} \in \mathring{C}\) . The map \(\lambda \mapsto \langle x_{0}, \xi(\lambda) \rangle\) tends to infinity when \(\lambda\) tends to \(\lambda_{m}\) . Prove that the set of values \(\xi(\lambda)/\langle x_{0}, \xi(\lambda) \rangle\) has a cluster value \(\xi_{m}\) such that:
\end{enumerate}

\begin{enumerate}
\def\labelenumi{(\roman{enumi})}
\item
  \(\xi_{\mathrm{m}}\in \mathrm{C}^{\circ}\) ;
\item
  \(\langle x_{0},\xi_{m}\rangle=1;\)
\item
  \(\xi_{\mathrm{m}} = \lambda_{\mathrm{m}}^{t}u(\xi_{\mathrm{m}})\) .
\end{enumerate}

\begin{enumerate}
\def\labelenumi{\alph{enumi})}
\setcounter{enumi}{2}
\item
  Let \(f \in E\) such that \(\langle f, \xi_{m} \rangle = 0\) . For \(\lambda \in R_{+}\) such that \(\lambda < \lambda_{m}\) , we denote by \(x(\lambda)\) the unique solution \(x\) in E of the equation \(x - \lambda u(x) = f\) . Show that the set of \(x(\lambda)\) is bounded in E (let d be a metric defining the topology of E; if there existed a sequence \(\lambda_{n} \to \lambda_{m}\) such that \(x(\lambda_{n})\) tends to infinity, set \(\hat{x}_{n} = x(\lambda_{n}) / d(0, x(\lambda_{n}))\) , and show that the sequence \((\hat{x}_{n})\) converges to an element \(\hat{x}\) such that \(d(0, \hat{x}) = 1\) , \(\hat{x} = \lambda u(\hat{x})\) and \(\langle \hat{x}, \xi_{m} \rangle = 0\) , which contradicts the fact that one must have \(\hat{x} \in \mathring{C}\) ).
\item
  There exists a unique element \(x \in E\) such that \(f = x - \lambda_{\mathrm{m}} u(x)\) and \(\langle x, \xi_{\mathrm{m}} \rangle = 0\) . (For existence, show using the preceding question that there exists a sequence \((\lambda_{n})\) converging to \(\lambda_{m}\) such that \((x(\lambda_{n}))\) converges to an element x satisfying the required conditions.)
\item
  The endomorphism \(1_{\mathrm{E}'} - \lambda_{\mathrm{m}}{}^{t}u\) of \(\mathrm{E}'\) induces, by passage to subspaces, a bijection of the polar \(x_{\mathrm{m}}^{\circ}\) of \(x_{\mathrm{m}}\) .
\item
  The eigenspace \(\operatorname{Ker}(1_{\mathrm{E}} - \lambda_{\mathrm{m}}u)\) is of dimension 1 and is equal to the polar of the image of \(1_{E} - \lambda_{m}u\) .
\end{enumerate}

CHAPTER IV

\chapter{Hilbertian spectral theory}
\setcounter{exercisepar}{0}

For every vector space E, we denote by \(1_{E}\) the identity map of E. If E, F and G are vector spaces and if v: \(E \rightarrow F\) and u: \(F \rightarrow G\) are linear maps, the linear map \(u \circ v\) from E to G will sometimes be denoted uv.

Let E be a pre-Hilbert space. We denote by \(\langle x|y\rangle_{E}\) , or simply \(\langle x|y\rangle\) , the inner product on E. The orthogonal of a subset A of E is denoted A°. It is a closed subspace of E; it is reduced to 0 if and only if A is total in E (EVT, V, p. 16, cor. 1). If F ⊂ E is a vector subspace of E, then F°° = \(\overline{F}\) (cf. EVT, V, p. 13). The spectrum of an endomorphism of E is always relative to the Banach algebra \(\mathcal{L}(E)\) .

Let F be a Hilbert space. We denote by \(\mathcal{L}_1(\mathrm{E})\) the space of endomorphisms of finite trace of E (EVT, V, p. 50, def. 8) and \(\mathcal{L}_2(\mathrm{E};\mathrm{F})\) the space of Hilbert-Schmidt maps from E to F (EVT, V, p. 51, def. 9). The map \(u\mapsto \| u\| _2 = (\mathrm{Tr}(u^{*}u))^{1 / 2}\) is a Hilbertian norm on \(\mathcal{L}_2(\mathrm{E};\mathrm{F})\) (EVT, V, p. 52, th. 1, (i)).

The space \(\mathcal{L}^{\mathrm{f}}(\mathrm{E};\mathrm{F})\) of continuous linear maps of finite rank is dense in \(\mathcal{L}^{\mathrm{c}}(\mathrm{E};\mathrm{F})\) (prop. 10 of III, p. 15 and cor. of prop. 11 of III, p. 16). In particular, the space \(\mathcal{L}_2(\mathrm{E};\mathrm{F})\) is dense in \(\mathcal{L}^{\mathrm{c}}(\mathrm{E};\mathrm{F})\) .

The Hilbertian tensor product of E and F (EVT, V, p. 28, def. 1) is denoted E \(\widehat{\otimes}_2\) F.

If X is a locally compact topological space and \(\mu\) a measure on X, we denote by \(\mathcal{L}^{p}(\mathrm{X},\mu)=\mathcal{L}_{\mathbf{C}}^{p}(\mathrm{X},\mu)\) and \(\mathrm{L}^{p}(\mathrm{X},\mu)=\mathrm{L}_{\mathbf{C}}^{p}(\mathrm{X},\mu)\) for \(p\in[1,+\infty]\) .

\section{COMPACT OPERATORS ON A HILBERT SPACE}

In this section, K is a field equal to R or C and E denotes a Hilbert space over the field K.

\subsection{Diagonal endomorphisms}

DEFINITION 1. --- Let \(\mathrm{B} = (e_i)_{i \in \mathrm{I}}\) be an orthonormal basis of E. An endomorphism \(u\) is said to be diagonal in the basis B or diagonal with respect to B if there exists a family \(\lambda = (\lambda_i)_{i \in \mathrm{I}}\) of elements of K such that \(u(e_i) = \lambda_i e_i\) for all \(i \in \mathrm{I}\).

Let \(\mathrm{B} = (e_i)_{i \in \mathrm{I}}\) be an orthonormal basis of E. We denote by \(\mathcal{D}_{\mathrm{B}}(\mathrm{E})\) the set of endomorphisms of E that are diagonal with respect to B. This is a closed, unital commutative subalgebra of \(\mathcal{L}(\mathrm{E})\) . Let \(u \in \mathcal{D}_{\mathrm{B}}(\mathrm{E})\) . The family \(\lambda = (\lambda_i)_{i \in \mathrm{I}}\) such that \(u(e_i) = \lambda_i e_i\) is determined uniquely by u and by the basis B, and it is called the family of eigenvalues of u relative to B.

Suppose that K = R and identify E with a subspace of \(\mathrm{E}_{(\mathbf{C})}\) . Let \(\mathrm{B} = (e_i)_{i \in \mathrm{I}}\) be an orthonormal basis of E. Denote by \(\mathrm{B}_{(\mathbf{C})}\) the orthonormal basis \((1 \otimes e_i)_{i \in \mathrm{I}}\) of \(\mathrm{E}_{(\mathbf{C})}\) (cf. EVT, V, p. 29, cor. 1). The map \(u \mapsto u_{(\mathbf{C})}\) induces an injective morphism of algebras from \(\mathcal{D}_{\mathrm{B}}(\mathrm{E})\) into \(\mathcal{D}_{\mathrm{B}_{(\mathbf{C})}}(\mathrm{E}_{(\mathbf{C})})\) ; its image is the set of \(u \in \mathcal{D}_{\mathrm{B}_{(\mathbf{C})}}(\mathrm{E}_{(\mathbf{C})})\) such that \(u(\mathrm{E}) \subset \mathrm{E}\) ; in other words, the set of diagonal endomorphisms u in the basis \(\mathrm{B}_{(\mathbf{C})}\) whose eigenvalues are real.

For every \(i \in \mathrm{I}\) , we denote by \(p_i\) the orthogonal projector of E with image \(\mathrm{Ke}_i\) . We have \(\| p_i \| = 1\) for all \(i \in \mathrm{I}\) . The endomorphism \(p_i\) is diagonal in the basis B, and the family of its eigenvalues is the family \((\delta_{ij})_{j\in\mathrm{I}}\) (Kronecker symbol, cf. A, II, p. 24).

Let \(\mathrm{B}^{\prime}=(f_{i})_{i\in\mathrm{I}}\) be an orthonormal basis of E and \(u:E\to E\) the isometric isomorphism such that \(u(f_{i})=e_{i}\) . Then \(\mathcal{D}_{\mathrm{B}^{\prime}}(\mathrm{E})=u^{-1}\mathcal{D}_{\mathrm{B}}(\mathrm{E})u\) .

We endow I with the discrete topology and denote by \(\mathcal{B}(\mathrm{I}) = \mathcal{C}_{b}(\mathrm{I}; \mathrm{K})\) the unital Banach algebra of bounded functions on I with values in K (example 2 of I, p. 17); if K = C, it is a star algebra (example 2 of I, p. 102).

Lemma 1. --- Let \(u \in \mathcal{D}_{\mathrm{B}}(\mathrm{E})\) and let \(\lambda = (\lambda_i)_{i \in \mathrm{I}}\) be the family of its eigenvalues. The family \(\lambda\) is bounded and one has \(\sup_i |\lambda_i| = \| u\|\) .

One has \(|\lambda_i| \leqslant \|u\|\) for all \(i\) , whence \(\sup_{i \in I} |\lambda_i| \leqslant \|u\|\) . Since moreover

\[
\| u (x) \| ^ {2} = \sum_ {i \in \mathrm{I}} | \lambda_ {i} | ^ {2} | \langle e _ {i} | x \rangle | ^ {2} \leqslant \left(\sup _ {i \in \mathrm{I}} | \lambda_ {i} | ^ {2}\right) \| x \| ^ {2}
\]

for all \(x \in \mathrm{E}\) (EVT, V, p. 22, prop. 5), it follows that \(\| u\| = \sup |\lambda_i|\) .

PROPOSITION 1. --- a) The map \(\alpha\) from \(\mathcal{D}_{\mathrm{B}}(\mathrm{E})\) into \(\mathcal{B}(\mathrm{I})\) that associates with a diagonal endomorphism \(u\) the family of eigenvalues of \(u\) is an isometric isomorphism of Banach algebras over K. If K = C, it is a morphism of involutive algebras.

\begin{enumerate}
\def\labelenumi{\alph{enumi})}
\setcounter{enumi}{1}
\item
  For any bounded family \(\lambda = (\lambda_i)_{i\in \mathrm{I}}\) the family \((\lambda_ip_i)_{i\in \mathrm{I}}\) is summable in the space \(\mathcal{L}(\mathrm{E})\) endowed with the topology of pointwise convergence, and the sum of this family is the unique mapping \(u\in \mathcal{D}_{\mathrm{B}}(\mathrm{E})\) such that \(\alpha (u) = \lambda\) ;
\item
  Let \(u \in \mathcal{D}_{\mathrm{B}}(\mathrm{E})\) and let \(\lambda\) be the family of its eigenvalues relative to B. The spectrum of \(u\) is the closure in K of the set of values of \(\lambda\) ;
\item
  If K = C, then the continuous functional calculus map of \(\mathcal{C}(\mathrm{Sp}(u))\) in \(\mathcal{L}(\mathrm{E})\) associates with a continuous function \(f \in \mathcal{C}(\mathrm{Sp}(u))\) the endomorphism \(\alpha(f \circ \lambda)\) in \(\mathcal{D}_{\mathrm{B}}(\mathrm{E})\) .
\end{enumerate}

By Lemma 1, the family of eigenvalues of an endomorphism diagonal in B is bounded, and the map from \(\mathcal{D}_{\mathrm{B}}(\mathrm{E})\) into \(\mathcal{B}(\mathrm{I})\) thus defined is a continuous isometric morphism of unital Banach algebras.

Let \(i \in I\) . For every \(j \in I\) , we have \(\langle u^{*}(e_{i}) | e_{j} \rangle = \lambda_{j} \langle e_{i} | e_{j} \rangle = \langle \overline{\lambda_{j}} e_{i} | e_{j} \rangle\) . It follows that \(u^{*}(e_{i}) = \overline{\lambda}_{i} e_{i}\) . The adjoint of u is therefore the diagonal endomorphism in the basis B for which \((\overline{\lambda}_{i})_{i \in I}\) is the family of eigenvalues. Assertion a) follows from this.

Let \(\lambda = (\lambda_i)_{i\in \mathrm{I}}\in \mathcal{B}(\mathrm{I})\) . For every \(x\in \mathrm{E}\) the family \((|\langle e_i|x\rangle |^2)_{i\in \mathrm{I}}\) is summable, with sum \(\| x\|^2\) (EVT, V, p. 22, prop. 5), whence, for every finite subset J of I:

\[
\begin{array}{r l r} & & {\left\| \sum_ {i \in \mathrm{J}} \lambda_ {i} p _ {i} (x) \right\| ^ {2} = \sum_ {i \in \mathrm{J}} | \lambda_ {i} | ^ {2} | \langle e _ {i} | x \rangle | ^ {2} \leqslant \Bigl (\underset {i \in \mathrm{I}} {\sup} | \lambda_ {i} | ^ {2} \Bigr) \sum_ {i \in \mathrm{J}} | \langle e _ {i} | x \rangle | ^ {2}} \\ & & {\leqslant \Bigl (\underset {i \in \mathrm{I}} {\sup} | \lambda_ {i} | \Bigr) ^ {2} \| x \| ^ {2}.} \end{array}
\]

Consequently, the family \((\lambda_{i}p_{i})\) is summable in \(\mathcal{L}(\mathrm{E})\) endowed with the topology of simple convergence. Its sum \(u_{\lambda}\) satisfies \(u_{\lambda}(e_{i}) = \lambda_{i}e_{i}\) ; it is therefore a diagonal endomorphism in the basis B, with eigenvalues \(\lambda\) . Assertion b) follows from this.

The last assertions follow from \(a\) ), from example 3 of I, p. 17 and from example 4 of I, p. 111.

Remark. --- The Banach algebra \(\mathcal{D}_{\mathrm{B}}(\mathrm{E})\) is a maximal commutative closed subalgebra of \(\mathcal{L}(\mathrm{E})\) . Indeed, let u be an endomorphism of E that commutes with \(\mathcal{D}_{\mathrm{B}}(\mathrm{E})\) . Let \(i \in I\) . Since the orthogonal projector \(p_{i}\) is diagonal in the basis B, we have \(p_{i}(u(e_{i})) = u(p_{i}(e_{i})) = u(e_{i})\) , which implies that \(u(e_{i})\) is proportional to \(e_{i}\) . Thus u is diagonal in the basis B.

If E is infinite-dimensional, there exist maximal commutative involutive subalgebras of \(\mathcal{L}(\mathrm{E})\) that are not isomorphic to \(\mathcal{D}_{\mathrm{B}}(\mathrm{E})\) (exercise 5 of IV, p. 314).

PROPOSITION 2. --- Let \(u \in \mathcal{D}_{\mathrm{B}}(\mathrm{E})\) and \(\lambda = (\lambda_i)_{i \in \mathrm{I}}\) be the family of its eigenvalues.

\begin{enumerate}
\def\labelenumi{\alph{enumi})}
\tightlist
\item
  The following conditions are equivalent:
\end{enumerate}

\begin{enumerate}
\def\labelenumi{(\roman{enumi})}
\item
  One has \(\lambda \in \mathcal{C}_{0}(\mathrm{I};\mathrm{K})\) ;
\item
  The endomorphism u is compact;
\item
  The family \((\lambda_{i}p_{i})_{i\in\mathrm{I}}\) is summable in the Banach space \(\mathcal{L}(\mathrm{E})\) . Its sum is then equal to u.
\end{enumerate}

\begin{enumerate}
\def\labelenumi{\alph{enumi})}
\setcounter{enumi}{1}
\tightlist
\item
  Suppose that \(u\) is compact and let \(\Lambda\) denote the set of values of \(\lambda\) . The set of \(i \in I\) such that \(\lambda_i \neq 0\) is countable. If E is infinite-dimensional, we have \(\mathrm{Sp}_{\mathrm{s}}(u) = \Lambda - \{0\}\) and \(\mathrm{Sp}(u) = \Lambda \cup \{0\}\) . If E is finite-dimensional, then \(\mathrm{Sp}_{\mathrm{s}}(u) = \mathrm{Sp}(u) = \Lambda\) .
\end{enumerate}

By Lemma 1, we have \(\| \sum_{j\in J}\lambda_jp_j\| = \sup_{j\in J}|\lambda_j|\) for every finite subset J of I. It follows, by the Cauchy criterion, that the family \((\lambda_ip_i)\)

is summable in \(\mathcal{L}(\mathrm{E})\) if and only if the family \(\lambda\) tends to 0 at infinity, which implies that conditions (i) and (iii) are equivalent.

Condition (iii) implies that \(u\) is compact (corollary of proposition 2 of III, p. 4). Conversely, suppose that the endomorphism \(u\) is compact. Since the family B is bounded in E, its image under \(u\) in E is relatively compact in E (III, p. 2), hence precompact. Let \(\varepsilon > 0\) and let \(J\) be a finite subset of I such that the image of B under \(u\) is contained in the union of balls of radius \(\varepsilon\) and centers \(u(e_j)\) for \(j \in \mathrm{J}\) . Let \(i \in \mathrm{I} - \mathrm{J}\) . There exists \(j \in \mathrm{J}\) such that \(\| u(e_i) - u(e_j) \| \leqslant \varepsilon\) , hence

\[
\left| \lambda_ {i} \right| ^ {2} \leqslant \left| \lambda_ {i} \right| ^ {2} + \left| \lambda_ {j} \right| ^ {2} = \left\| u \left(e _ {i}\right) - u \left(e _ {j}\right) \right\| ^ {2} \leqslant \varepsilon^ {2}.
\]

Consequently, we have \(\lambda \in \mathcal{C}_0(\mathrm{I};\mathrm{K})\) , that is, condition (i).

Finally, the spectrum of u and its sensitive spectrum are computed in terms of \(\lambda\) using prop. 1, c) and proposition 5 of III, p. 90.

Remark. --- When I is infinite, condition \(\lambda\in\mathcal{C}_{0}(\mathrm{I};\mathrm{K})\) can also be stated as “the family \(\lambda\) tends to 0 along the filter of complements of finite subsets of I.”

\subsection{Diagonalization of compact endomorphisms}

THEOREM 1. --- Suppose K = C. Let u be a compact and normal endomorphism of E. There exists an orthonormal basis B of E such that u is diagonal in the basis B.

The set \(\mathrm{Sp}_{\mathrm{s}}(u)\) is countable, and it does not contain 0 if E is infinite-dimensional (III, p. 90, prop. 5, \(b\) ). For every element \(\lambda \in \mathrm{Sp}_{\mathrm{s}}(u)\) , denote by \(\mathrm{N}_{\lambda}\) the nilspace of \(u - \lambda 1_{\mathrm{E}}\) . It is finite-dimensional (loc. cit.) and, since \(u\) is normal, it coincides with the eigenspace of \(u\) relative to \(\lambda\) (EVT, V, p. 43, cor. of prop. 8). The spaces \(\mathrm{N}_{\lambda}\) are pairwise orthogonal (I, p. 132, n° 5).

Let F be the subspace of E that is the Hilbert sum of the spaces \(N_{\lambda}\) for \(\lambda \in \mathrm{Sp}_{\mathrm{s}}(u) - \{0\}\) . It is a space of countable type, stable under u since each subspace \(N_{\lambda}\) is stable under u. Since \(N_{\lambda}\) is also the eigenspace of \(u^{*}\) relative to \(\overline{\lambda}\) (EVT, V, loc. cit.), the endomorphism u induces an endomorphism \(\widetilde{u}\) of \(F^{\circ}\) by passing to subspaces. The endomorphism \(\widetilde{u}\) is compact (prop. 3 of III, p. 5) and normal (lemma 4 of I, p. 135). By construction, the sensitive spectrum of \(\widetilde{u}\) is contained in \(\{0\}\) (indeed, every eigenvector for \(\widetilde{u}\) would be one for u, hence would belong to one of the spaces \(N_{\lambda}\) if the corresponding eigenvalue were nonzero). Thus the spectral radius of \(\widetilde{u}\) is zero, whence \(\widetilde{u}=0\) since \(\widetilde{u}\) is normal (cor. 1 of I, p. 108). We therefore have \(F^{\circ}\subset\operatorname{Ker}(u)\) .

For every \(\lambda \in \mathrm{Sp}_{\mathrm{s}}(u)\) , let \(B_{\lambda}\) be an orthonormal basis of \(N_{\lambda}\) and let \((e_j)_{j \in J}\) be the family formed by the union of the \(B_{\lambda}\) ; it is an orthonormal basis of F. Let B be the union of \((e_j)_{j \in J}\) and of an orthonormal basis of \(F^{\circ}\) . This is an orthonormal basis of E and \(u\) is diagonal in the basis B.

COROLLARY 1. --- Let u be a compact Hermitian endomorphism of E. There exists an orthonormal basis B of E such that u is diagonal in the basis B and its eigenvalues are real.

If K = C, this follows immediately from the theorem. Suppose K = R. The space E is an R-structure on \(\mathrm{E}_{(\mathbf{C})}\) (cf. A, II, p. 119). The endomorphism \(u_{(\mathbf{C})}\) of \(\mathrm{E}_{(\mathbf{C})}\) is compact (remark 4 of III, p. 2) and Hermitian. Let \(B = (e_j)_{j \in J}\) be an orthonormal basis of \(\mathrm{E}_{(\mathbf{C})}\) such that \(u_{(\mathbf{C})} \in \mathcal{D}_{\mathrm{B}}(\mathrm{E}_{(\mathbf{C})})\) and let \(\lambda\) be the family of eigenvalues of \(u_{(\mathbf{C})}\) (theorem 1). We have \(\lambda \in R^{J}\) (I, p. 106, prop. 4); since the linear map \(u_{(\mathbf{C})}\) is R-rational, the eigenspace of \(u_{(\mathbf{C})}\) relative to \(\lambda_j\) is R-rational for every \(j \in J\) (cf. A, V, p. 60, prop. 6). Hence there exists a basis of it belonging to E, and a fortiori, there exists an orthonormal basis in E. The union of these bases is an orthonormal basis \(B_R\) of E such that \(u \in \mathcal{D}_{\mathrm{B}_R}(\mathrm{E})\) .

COROLLARY 2. --- Let F be a Hilbert space and let u be a compact continuous linear map from E to F. There exists a countable set I, an orthonormal basis \((e_{i})_{i\in\mathrm{I}}\) of the initial space \(\operatorname{Ker}(u)^{\circ}\) of u, an orthonormal family \((f_{i})_{i\in\mathrm{I}}\) of F, and a family \((\alpha_{i})_{i\in\mathrm{I}}\in(\mathbf{R}_{+}^{*})^{\mathrm{I}}\) such that \(u(e_{i})=\alpha_{i}f_{i}\) for all \(i\in I\) .

Let \(v = u^{*} \circ u\) . This is a compact, positive, hence Hermitian, endomorphism of E. By corollary 1, there exists an orthonormal basis \((e_{j})_{j \in \mathrm{J}}\) of E such that \(v\) is diagonal in this basis. The family \((\lambda_{j})_{j \in \mathrm{J}}\) of its eigenvalues is contained in \(\mathbf{R}_{+}^{\mathrm{J}}\) . Set \(\alpha_{j} = \sqrt{\lambda_{j}}\) for all \(j \in \mathrm{J}\) . Let I be the set of \(j \in \mathrm{J}\) such that \(\alpha_{j} \neq 0\) . This is a countable set since \(v\) is compact. The family \((e_{i})_{i \in \mathrm{I}}\) is an orthonormal basis of the initial space of \(v\) which is the initial space \(\operatorname{Ker}(u)^{\circ}\) of \(u\) (EVT, V, p. 43, prop. 8). Set \(f_{i} = \frac{1}{\alpha_{i}} u(e_{i})\) for \(i\in \mathrm{I}\) . For any \(i\) and \(j\) in I, we have

\[
\langle f _ {i} | f _ {j} \rangle = \frac {1}{\alpha_ {i} \alpha_ {j}} \langle u (e _ {i}) | u (e _ {j}) \rangle = \frac {1}{\alpha_ {i} \alpha_ {j}} \langle v (e _ {i}) | e _ {j} \rangle = \frac {\lambda_ {i}}{\alpha_ {i} \alpha_ {j}} \langle e _ {i} | e _ {j} \rangle ,
\]

whence it follows that the family \((f_{i})_{i\in\mathrm{I}}\) is orthonormal in F. The corollary follows, since \(u(e_{i})=\alpha_{i}f_{i}\) for all \(i\in I\) .

DEFINITION 2. --- With the notations of the corollary, the family \((\alpha_{i})_{i\in I}\) is the family of singular values of \(u\) relative to the orthonormal basis \((e_{i})_{i\in I}\) of the initial space of \(u\) .

Remarks. --- 1) This corollary generalizes th. 2 of EVT, V, p. 54, which corresponds to Hilbert--Schmidt maps.

\begin{enumerate}
\def\labelenumi{\arabic{enumi})}
\setcounter{enumi}{1}
\tightlist
\item
  With the notations of the corollary, we have the formula
\end{enumerate}

\[
u (x) = \sum_ {i \in I} \alpha_ {i} \langle e _ {i} | x \rangle f _ {i}\tag{1}
\]

for all \(x \in E\) .

\begin{enumerate}
\def\labelenumi{\arabic{enumi})}
\setcounter{enumi}{2}
\tightlist
\item
  Let \(u\) be a compact positive endomorphism of E. Let \(B = (f_i)_{i \in I}\) be an orthonormal basis of E such that \(u\) is diagonal in the basis B (theorem 1), and let \((\lambda_i)_{i \in I}\) be the family of eigenvalues of \(u\) in this basis. Let J be the set of \(i \in I\) such that \(\lambda_i > 0\) ; the family \((e_i)_{i \in J}\) is an orthonormal basis of the space \(\text{Ker}(u)^{\circ}\) . For every \(i \in J\) , set \(e_i = f_i\) and \(\alpha_i = \lambda_i\) . We obtain \(u(e_i) = \alpha_if_i\) : the family \((\alpha_i)_{i \in J}\) is the family of singular values of \(u\) relative to the basis \((e_i)_{i \in J}\) .
\end{enumerate}

\subsection{Decreasing sequence of eigenvalues}

In this section, we assume that K = C.

We set \(\overline{\mathbf{N}} = \mathbf{N} \cup \{+\infty\} \subset \overline{\mathbf{R}}\) . In this section, we say that a vector space E has dimension \(+\infty \in \overline{\mathbf{N}}\) if E is not finite-dimensional. Let \(\mathrm{I_E} \subset \mathbf{N}\) be the set of dimensions of finite-dimensional subspaces F of E such that \(\mathrm{F} \neq \mathrm{E}\) . We have \(\mathrm{I_E} = \mathbf{N}\) if E is infinite-dimensional, and otherwise \(\mathrm{I_E} = \{0, \ldots, \dim(\mathrm{E}) - 1\}\) . We set \(\mathrm{I} = \mathrm{I_E}\) when no confusion can result.

Let \(u\) be a compact positive (in particular Hermitian) endomorphism of E. The sensitive spectrum of \(u\) is the set of values of a strictly decreasing sequence \((\nu_{k})_{0\leqslant k<\mathrm{Card}(\mathrm{Sp}_{\mathrm{s}}(u))}\) of positive real numbers (cf.

prop. 5 of III, p. 90). For every integer \(k\) such that \(0 \leqslant k < \text{Card}(\text{Sp}_s(u))\) , we denote by \(n_k\) the spectral multiplicity of \(\nu_k\), with \(n_k \geqslant 1\) . Let \(\text{M} \in \overline{\mathbf{N}}\) be the sum of the spectral multiplicities \(n_k\) ; it is the dimension of the image of \(u\) . We have \(\text{M} \leqslant \text{Card}(\text{I})\) .

For \(0 \leqslant n < M\) , one defines \(\lambda_n(u) = \nu_k\) , where \(k \geqslant 0\) is the unique integer such that

\[
n _ {0} + \dots + n _ {k - 1} \leqslant n <   n _ {0} + \dots + n _ {k}.
\]

One sets \(\lambda_{n}(u)=0\) if \(n\in I\) satisfies \(n\geqslant M\) . This case can occur only if I=N and if \(\mathrm{Sp}_{\mathrm{s}}(u)\) is finite (or, equivalently, if E is infinite-dimensional and u is of finite rank).

The sequence \((\lambda_{n}(u))_{n\in\mathrm{I}}\) is decreasing; for every \(\lambda\in\mathrm{Sp}_{\mathrm{s}}(u)\) , the number of integers n such that \(\lambda_{n}(u)=\lambda\) is equal to the spectral multiplicity of the eigenvalue \(\lambda\) of u.

One says, by abuse of language, that \((\lambda_{n}(u))_{n\in\mathrm{J}}\) is the decreasing sequence of eigenvalues of u repeated with their multiplicities.

PROPOSITION 3. --- Let u be a compact positive endomorphism of E. There exists an orthonormal family \((e_{n})_{n\in\mathrm{I}}\) in E such that, for every \(x\in E\) , one has

\[
u (x) = \sum_ {n \in \mathrm{I}} \lambda_ {n} (u) \langle e _ {n} | x \rangle e _ {n}, \quad \langle x | u (x) \rangle = \sum_ {n \in \mathrm{I}} \lambda_ {n} (u) | \langle e _ {n} | x \rangle | ^ {2}.
\]

Let \(\mathrm{B} = (f_{j})_{j \in \mathrm{J}}\) be an orthonormal basis of E in which u is diagonal (cor. 1 of IV, p. 150) and \((\lambda_{j})_{j \in \mathrm{J}}\) the family of eigenvalues of u in the basis B. Let J' be the set of \(j \in J\) such that \(\lambda_{j}\) belongs to the sensitive spectrum of u.

For each \(\lambda \in \mathrm{Sp}_{\mathrm{s}}(u)\) , there exists a bijection between the integers \(n\) such that \(0 \leqslant n < \mathrm{M}\) and \(\lambda_{n}(u) = \lambda\) and the \(j \in \mathrm{J}'\) such that \(\lambda_{j} = \lambda\) , since these two sets have the same cardinality, namely the spectral multiplicity of \(\lambda\) . A choice of such bijections for every \(\lambda\) defines a bijection \(\iota\) from the set of integers such that \(0 \leqslant n < \mathrm{M}\) to the set \(\mathrm{J}'\) . One defines the sequence \((e_n)_{0 \leqslant n < \mathrm{M}}\) in E by setting \(e_n = f_{\iota(n)}\) for \(0 \leqslant n < \mathrm{M}\) . It is an orthonormal family in E.

In the case where M \textless{} Card(I) = +∞, the space F generated by \(\{e_{0},\ldots,e_{M-1}\}\) is finite-dimensional and its orthogonal F° is infinite-dimensional; one chooses for \((e_{n})_{n\geqslant M}\) an orthonormal family in F°.

For every \(x \in \mathrm{E}\) , we have

\[
u (x) = \sum_ {j \in J ^ {\prime}} \lambda_ {j} \langle f _ {j} \mid x \rangle f _ {j} = \sum_ {0 \leqslant n <   M} \lambda_ {n} (u) \langle e _ {n} \mid x \rangle e _ {n}
\]

since \(u(f_{j}) = 0\) when \(j \in J - J'\) . If \(n \in I\) satisfies \(n \geqslant M\) , we have \(\lambda_{n}(u) = 0\) , and one obtains the first formula of the proposition. The second follows from it.

PROPOSITION 4. --- Let u be a compact positive endomorphism of E. Let \(f \in \mathcal{C}(\mathbf{R}_{+})\) be a continuous increasing function such that \(f(0) = 0\) .

The endomorphism \(f(u)\) is compact and positive and, for every \(n \in \mathrm{I}_{\mathrm{E}}\) , we have \(\lambda_n(f(u)) = f(\lambda_n(u))\) .

The endomorphism \(f(u)\) is compact and positive by prop. 6, b) of III, p. 91 and prop. 15, a) of I, p. 117. The spectrum of \(f(u)\) is the image under \(f\) of the spectrum of \(u\) (cor. 2 of I, p. 111). If \(\lambda \in \mathrm{Sp}_{\mathrm{s}}(f(u))\) , the spectral multiplicity of \(\lambda\) is the sum of the spectral multiplicities of the \(\mu \in \mathrm{Sp}(u)\) such that \(f(\mu) = \lambda\) (cor. 2 of III, p. 84). Since \(f\) is increasing, the sequence \((f(\lambda_n(u)))_{n \in \mathrm{I_E}}\) is decreasing. The assertion then follows from the definition of the sequence \((\lambda_n(f(u)))_{n \in \mathrm{I_E}}\) .

\subsection{Variational characterizations of eigenvalues}

In this section, we assume that K = C.

Let \(u\) be a compact and positive endomorphism of E and let \((\lambda_n(u))_{n\in \mathrm{I_E}}\) be the decreasing sequence of eigenvalues of \(u\) . We set \(\mathrm{I} = \mathrm{I_E}\) .

For every closed subspace F of E, we denote by

\[
r _ {\mathrm{F}} (u) = \inf _ {x \in \mathrm{F} - \{0 \}} \frac {\langle x \mid u (x) \rangle}{\| x \| ^ {2}}, \quad \mathrm{R} _ {\mathrm{F}} (u) = \sup _ {x \in \mathrm{F} ^ {\circ} - \{0 \}} \frac {\langle x \mid u (x) \rangle}{\| x \| ^ {2}},
\]

where the lower bound (resp. the upper bound) is taken in \([0, +\infty]\) .

For every \(n \in \mathbf{N}\) , we denote by \(\mathcal{F}_n\) the set of vector subspaces \(\mathrm{F} \subset \mathrm{E}\) of dimension \(n\) . We say that a subspace \(\mathrm{F} \in \mathcal{F}_n\) is adapted to \(u\) if it admits an orthonormal basis \((f_i)_{0 \leqslant i \leqslant n-1}\) such that \(u(f_i) = \lambda_i(u)f_i\) for \(0 \leqslant i \leqslant n-1\) .

PROPOSITION 5. --- Let \(n \in I\) .

\begin{enumerate}
\def\labelenumi{\alph{enumi})}
\item
  For every subspace \(\mathrm{F} \in \mathcal{F}_{n+1}\) adapted to \(u\) , we have \(\lambda_n(u) = r_{\mathrm{F}}(u)\) ;
\item
  For every subspace \(F \in F_{n}\) adapted to u, we have \(\lambda_{n}(u) = R_{\mathrm{F}}(u)\) .
\end{enumerate}

Let \(F \in F_{n+1}\) be a subspace adapted to u, and \((f_i)_{0 \leqslant i \leqslant n}\) be an orthonormal basis of F such that \(u(f_i) = \lambda_i(u)f_i\) for all i. For all x in F, we have

\[
\langle x \mid u (x) \rangle = \sum_ {0 \leqslant i \leqslant n} \lambda_ {i} (u) | \langle f _ {i} \mid x \rangle | ^ {2} \geqslant \lambda_ {n} (u) \sum_ {0 \leqslant i \leqslant n} | \langle f _ {i} \mid x \rangle | ^ {2} = \lambda_ {n} (u) \| x \| ^ {2},
\]

with equality if \(x = f_{n}\) . This implies that \(r_{\mathrm{F}}(u) = \lambda_{n}(u)\) whence assertion (a).

Let \(F \in F_{n}\) be a subspace adapted to u. The closed subspace \(F^{\circ}\) is nonzero (since \(n = \dim(F) < \dim(E)\) ) and stable under u; the endomorphism \(\widetilde{u}\) of \(F^{\circ}\) deduced from u by passing to subspaces is compact and positive.

One has \(\mathrm{Sp}(\widetilde{u})\subset \mathrm{Sp}(u)\) . We also have \(\mathrm{Sp}(\widetilde{u})\subset [0,\lambda_n(u)]\) . Indeed, it suffices to verify that \(\lambda \leqslant \lambda_{n}(u)\) for all \(\lambda \in \mathrm{Sp}_{\mathrm{s}}(\widetilde{u})\) . The number \(\lambda\) is then an eigenvalue of \(u\) , so there exists \(j\in I\) such that \(\lambda = \lambda_j(u)\) . The eigenspace of \(u\) relative to \(\lambda_{j}(u)\) is therefore not contained in F, which implies that \(\lambda_{j}(u)\leqslant \lambda_{n}(u)\) .

Suppose that \(\lambda_n(u) > 0\) . The eigenspace of \(u\) relative to \(\lambda_n(u)\) is then not contained in F, and \(\lambda_n(u)\) therefore belongs to the sensitive spectrum of \(\widetilde{u}\) . We deduce that \(\sup (\mathrm{Sp}(\widetilde{u})) = \lambda_n(u)\) . If \(\lambda_n(u) = 0\) , we have the same result since the spectrum of \(\widetilde{u}\) is then reduced to \(\{0\}\) .

Moreover, we have by definition \(\mathrm{R}_{\mathrm{F}}(u)=\mathrm{R}_{\{0\}}(\widetilde{u})\) , and finally \(\mathrm{R}_{\{0\}}(\widetilde{u})=\sup(\mathrm{Sp}(\widetilde{u}))\) according to prop. 9 of I, p. 139, a). Assertion b) is proved.

PROPOSITION 6. --- For all \(n \in I\) , we have

\[
\lambda_ {n} (u) = \sup _ {\mathrm{F} \in \mathscr {F} _ {n + 1}} r _ {\mathrm{F}} (u) = \inf _ {\mathrm{F} \in \mathscr {F} _ {n}} \mathrm{R} _ {\mathrm{F}} (u).
\]

Let \((e_{n})_{n\in\mathrm{I}}\) be an orthonormal family having the property of prop. 3 of IV, p. 152. For every integer \(n\) such that \(1\leqslant n<M+1\) , let \(F_{n}\in\mathcal{F}_{n}\) be the n-dimensional subspace of E generated by \((e_{0},\ldots,e_{n-1})\) ; by construction, the space \(F_{n}\) is adapted to u. According to prop. 5, we therefore have

\[
\lambda_ {n} (u) = r _ {\mathrm{F} _ {n + 1}} (u) = \mathrm{R} _ {\mathrm{F} _ {n}} (u).\tag{2}
\]

Let \(n \in I\) . Let \(F \in \mathcal{F}_{n+1}\) . The restriction to \(F\) of the orthogonal projector onto \(F_n\) is not injective, so there exists \(x \neq 0\) in \(F\) orthogonal to \(F_n\) .

Since \(x \in \mathrm{F}_{n}^{\circ}\) , we then have (prop. 3 of IV, p. 152)

\[
\begin{array}{l}\langle x\mid u(x)\rangle = \sum_{\substack{m\in \mathrm{I}\\ m\geqslant n}}\lambda_{m}(u)|\langle e_{m}\mid x\rangle |^{2}\\ \leqslant \lambda_{n}(u)\sum_{\substack{m\in \mathrm{I}\\ m\geqslant n}}|\langle e_{m}\mid x\rangle |^{2} = \lambda_{n}(u)\| x\|^{2}. \end{array}
\]

This proves that \(r_{\mathrm{F}}(u)\leqslant\lambda_{n}(u)\) , whence in particular the inequality

\[
\sup _ {\mathrm{F} \in \mathscr {F} _ {n + 1}} r _ {\mathrm{F}} (u) \leqslant \lambda_ {n} (u).\tag{3}
\]

Let \(F \in F_{n}\) . The restriction to \(F_{n+1}\) of the orthogonal projector onto F is not injective, so there exists a vector \(x \neq 0\) in \(F_{n+1}\) orthogonal to F. Since \(x \in F_{n+1}\) , we have (loc. cit.)

\[
\begin{array}{c} \langle x \mid u (x) \rangle = \sum_ {0 \leqslant m \leqslant n} \lambda_ {m} (u) | \langle e _ {m} \mid x \rangle | ^ {2} \\ \geqslant \lambda_ {n} (u) \sum_ {0 \leqslant m \leqslant n} | \langle e _ {m} \mid x \rangle | ^ {2} = \lambda_ {n} (u) \| x \| ^ {2} \end{array}
\]

and therefore \(\mathrm{R_F}(u) \geqslant \lambda_n(u)\) . In particular, it follows that

\[
\inf _ {\mathrm{F} \in \mathcal {F} _ {n}} \mathrm{R} _ {\mathrm{F}} (u) \geqslant \lambda_ {n} (u).\tag{4}
\]

In view of formulas (2), (3) and (4), the proposition is proved.

\subsection{Applications of the variational characterization of eigenvalues}

In this section, we consider Hilbert spaces over C.

PROPOSITION 7. --- Let u and v be positive compact endomorphisms of E.

\begin{enumerate}
\def\labelenumi{\alph{enumi})}
\item
  One has \(|\lambda_n(u) - \lambda_n(v)|\leqslant \| u - v\|\) for all \(n\in \mathbf{I}\) ;
\item
  If \(u \leqslant v\) , then \(\lambda_n(u) \leqslant \lambda_n(v)\) for all \(n \in I\) .
\end{enumerate}

Let \(n \in I\) and let F be an n-dimensional vector subspace of E. For all \(x \in F\) , we have

\[
| \langle x | v (x) \rangle - \langle x | u (x) \rangle | \leqslant \| u - v \| \| x \| ^ {2},
\]

therefore the inequalities

\[
\mathrm{R} _ {\mathrm{F}} (v) - \| u - v \| \leqslant \mathrm{R} _ {\mathrm{F}} (u) \leqslant \mathrm{R} _ {\mathrm{F}} (v) + \| u - v \|.
\]

The assertion \(a\) ) follows from this according to proposition 6 of IV, p. 154.

If \(u \leqslant v\) , we have \(\langle x | u(x) \rangle \leqslant \langle x | v(x) \rangle\) for all \(x \in E\) . For every \(n \in I\) and every subspace \(F\) of dimension \(n\) , we therefore have \(R_F(u) \leqslant R_F(v)\) , whence \(\lambda_n(u) \leqslant \lambda_n(v)\) (loc. cit.).

PROPOSITION 8. --- Let u be a positive compact endomorphism of E. Let H be a closed subspace of E and \(i_{\mathrm{H}}: \mathrm{H} \to \mathrm{E}\) the canonical injection. Denote by \(u_{\mathrm{H}}\) the endomorphism \(i_{\mathrm{H}}^{*}ui_{\mathrm{H}}\) of H. It is compact and positive.

\begin{enumerate}
\def\labelenumi{\alph{enumi})}
\item
  One has \(\mathrm{I_H}\subset \mathrm{I_E}\) and \(\lambda_n(u_{\mathrm{H}})\leqslant \lambda_n(u)\) for all \(n\in \mathrm{I_H}\) ;
\item
  If H has finite codimension \(k \in N\) in E, then we have \(I_{H} + k \subset I_{E}\) and \(\lambda_{n+k}(u) \leqslant \lambda_{n}(u_{\mathrm{H}})\) for all \(n \in I_{H}\) .
\end{enumerate}

The endomorphism \(u_{\mathrm{H}}\) is compact (prop. 3 of III, p. 5). It is positive because \(\langle x \mid u_{\mathrm{H}}(x) \rangle = \langle i_{\mathrm{H}}(x) \mid u(i_{\mathrm{H}}(x)) \rangle \geqslant 0\) for all \(x \in \mathrm{H}\) .

Let \(n \in I_{H} \subset I_{E}\) . Let F be a subspace of dimension \(n + 1\) of H adapted to \(u_{H}\) . We therefore have \(\lambda_{n}(u_{\mathrm{H}}) = r_{\mathrm{F}}(u_{\mathrm{H}})\) (prop 5 of IV, p. 153, a)), and since moreover \(r_{\mathrm{F}}(u_{\mathrm{H}}) = r_{\mathrm{F}}(u) \leqslant \lambda_{n}(u)\) (prop. 6 of IV, p. 154), we obtain assertion a).

Suppose that H has codimension \(k \in N\) in E and that \(n \in I_{H}\) . Let F be a subspace of H of dimension n adapted to \(u_{H}\) . Its orthogonal in H is equal to \(H \cap F^{\circ}\) , and it is the orthogonal in E of the subspace \(F + H^{\circ}\) of dimension \(n + k\) . Hence \(n + k \in I_{E}\) and (prop 5 of IV, p. 153, b))

\[
\lambda_{n}(u_{\mathrm{H}}) = \sup_{\substack{x\in \mathrm{H}\cap \mathrm{F}^{\circ}\\ x\neq 0}}\frac{\langle x\mid u_{\mathrm{H}}(x)\rangle}{\|x\|^{2}} = \mathrm{R}_{\mathrm{F} + \mathrm{H}^{\circ}}(u)
\]

whence \(\lambda_{n}(u_{\mathrm{H}})\leqslant \lambda_{n + k}(u)\) (prop. 6 of IV, p. 154).

DEFINITION 3. --- Let F be a Hilbert space and let u be a compact linear map from E into F.

For every integer \(n \in \mathrm{I}_{\mathrm{E}}\) we set \(\alpha_{n}(u) = \lambda_{n}(u^{*} \circ u)\) . Let J be the set of \(n \in \mathrm{I}_{\mathrm{E}}\) such that \(\alpha_{n}(u) > 0\) . The family \((\alpha_{n}(u))_{n \in \mathrm{J}}\) is called the sequence of singular values of \(u\) repeated with multiplicity.

We say that the sequence \((\alpha_{n}(u))_{n\in\mathrm{I}_{\mathrm{E}}}\) is the extended sequence of singular values of u.

The sequence \((\alpha_{n}(u))_{n\in\mathrm{I}_{\mathrm{E}}}\) is well defined since the endomorphism \(u^{*}\circ u\) of E is compact (III, p. 5, prop. 3); it is a decreasing family of positive real numbers since \(u^{*}\circ u\) is positive.

PROPOSITION 9. --- Let F be a Hilbert space and let u be a compact linear map from E into F.

\begin{enumerate}
\def\labelenumi{\alph{enumi})}
\tightlist
\item
  For \(n \in \mathrm{I}_{\mathrm{E}}\) , we have
\end{enumerate}

\[
\alpha_ {n} (u) = \sup _ {\mathrm{F} \in \mathscr {F} _ {n + 1}} \inf _ {x \in \mathrm{F} - \{0 \}} \frac {\| u (x) \|}{\| x \|} = \inf _ {\mathrm{F} \in \mathscr {F} _ {n}} \sup _ {x \in \mathrm{F} ^ {\circ} - \{0 \}} \frac {\| u (x) \|}{\| x \|}.
\]

\begin{enumerate}
\def\labelenumi{\alph{enumi})}
\setcounter{enumi}{1}
\tightlist
\item
  Let J be the set of \(n \in I_{E}\) such that \(\alpha_{n}(u) \neq 0\) . There exist orthonormal families \((e_{n})_{n \in \mathrm{J}}\) in E and \((f_{n})_{n \in \mathrm{J}}\) in F such that for every \(x \in E\) , we have
\end{enumerate}

\[
u (x) = \sum_ {n \in J} \alpha_ {n} (u) \langle e _ {n} | x \rangle f _ {n}.
\]

Since \(\langle x \mid u^{*}u(x)\rangle = \|u(x)\|^{2}\) for all \(x \in E\) , the definitions and prop. 6 of IV, p. 154, a), imply the equality of the first assertion. Let \((e_{n})_{n\in\mathrm{I}_{\mathrm{E}}}\) be an orthonormal family of E satisfying the conclusions of prop. 3 of IV, p. 152 applied to the compact positive endomorphism \(u^{*} \circ u\) of E. Set \(f_{n} = \alpha_{n}^{-1}u(e_{n})\) for \(n \in J\) . Reasoning as in the proof of corollary 2 of IV, p. 150, we obtain assertion c).

COROLLARY. --- Let F be a Hilbert space and u a compact linear map from E into F. Let \(w \in \mathcal{L}(\mathrm{E})\) and \(v \in \mathcal{L}(\mathrm{F})\) . For every \(n \in I_{E}\) , we have \(\alpha_{n}(w \circ u \circ v) \leqslant \|v\|\) \(\|w\|\) \(\alpha_{n}(u)\) .

This is a consequence of assertion \(a\) ) of the preceding proposition.

Remarks. --- 1) The sequence \((\alpha_{n}(u))_{n\in\mathrm{I}_{\mathrm{E}}}\) being decreasing, the set J is either equal to \(\mathrm{I}_{\mathrm{E}}\) , or equal to a segment \(\{0,\ldots,m\}\) in \(\mathrm{I}_{\mathrm{E}}\) where \(m\in\mathbf{N}\) . This latter case holds if and only if \(u\) has finite rank.

\begin{enumerate}
\def\labelenumi{\arabic{enumi})}
\setcounter{enumi}{1}
\item
  Since \(\| u(x)\| = \| |u|(x)\|\) for all \(x\in \mathrm{E}\) (I, p. 139, prop. 10), we have \(\alpha_{n}(u) = \alpha_{n}(|u|)\) for all \(n\in \mathrm{I}_{\mathrm{E}}\) .
\item
  If \(u\) is positive, then \(\alpha_{n}(u)=\lambda_{n}(u)\) for all \(n\in\mathrm{I}_{\mathrm{E}}\) (indeed, we then have \(\alpha_{n}(u)^{2}=\lambda_{n}(u^{*}u)=\lambda_{n}(u^{2})=\lambda_{n}(u)^{2}\) according to prop. 4 of IV, p. 153).
\item
  It is possible that \(\alpha_{n}(v\circ u\circ w)=0\) even if \(\alpha_{n}(u)\) is nonzero; in this case, \(\alpha_{n}(v\circ u\circ w)\) is not a singular value of \(v\circ u\circ w\) .
\end{enumerate}